% !TEX TS-program = pdflatex
\pdfoutput=1
\documentclass[11pt,a4paper]{article}

% ============================================================
% Packages
% ============================================================
\usepackage{amsmath,amsfonts,amssymb,amsthm,fullpage,booktabs}
\usepackage{float} % enables the [H] float placement specifier
\usepackage{algorithm,algpseudocode}
\usepackage[utf8]{inputenc}
\usepackage[numbers,sort&compress]{natbib}
\usepackage{hyperref}
\usepackage{xcolor}
\usepackage{enumitem}
\usepackage{microtype}
\usepackage{graphicx}
% Safe figure include: compile even if figure files are missing
\makeatletter
\newcommand{\safeincludegraphics}[2][]{%
	\begingroup
	% Make \IfFileExists search the same paths as \includegraphics
	\let\input@path\Ginput@path
	\def\@foundfile{}%
	% Try without extension first (rare), then common extensions
	\IfFileExists{#2}{\def\@foundfile{#2}}{%
		\IfFileExists{#2.pdf}{\def\@foundfile{#2.pdf}}{%
			\IfFileExists{#2.png}{\def\@foundfile{#2.png}}{%
				\IfFileExists{#2.jpg}{\def\@foundfile{#2.jpg}}{%
					\IfFileExists{#2.jpeg}{\def\@foundfile{#2.jpeg}}{%
						\IfFileExists{#2.eps}{\def\@foundfile{#2.eps}}{}%
					}%
				}%
			}%
		}%
	}%
	\ifx\@foundfile\@empty
	\fbox{\parbox[c][0.18\textheight][c]{0.9\linewidth}{\centering
			Missing figure: \texttt{\detokenize{#2}}}}%
	\else
	\includegraphics[#1]{\@foundfile}%
	\fi
	\endgroup
}
\makeatother

\usepackage{multirow}
\usepackage{subcaption}
\graphicspath{
	{./results/fig_multiobj/}
	{./results/fig_single/}
	{./results_scalability/fig_scalability/}
	{./results_robust/fig_robust/}
	{./results_param/param_study/}
}

\emergencystretch=2em
\allowdisplaybreaks

\hypersetup{
	colorlinks=true,
	linkcolor=blue,
	filecolor=magenta,
	urlcolor=cyan,
	citecolor=blue,
}

% Avoid hyperref PDF string warnings for math in bookmarks
\pdfstringdefDisableCommands{%
	\def\cO{O}%
	\def\one{1}%
	\def\J{J}%
	\def\Pmat{P}%
	\def\norm#1{||#1||}%
	\def\ip#1#2{<#1,#2>}%
}

% ============================================================
% Theorem-like
% ============================================================
\newtheorem{theorem}{Theorem}[section]
\newtheorem{lemma}[theorem]{Lemma}
\newtheorem{proposition}[theorem]{Proposition}
\newtheorem{assumption}[theorem]{Assumption}
\newtheorem{remark}[theorem]{Remark}

% ============================================================
% Macros
% ============================================================
\newcommand{\R}{\mathbb{R}}
\newcommand{\one}{\mathbf{1}}
\newcommand{\ip}[2]{\left\langle #1,#2\right\rangle}
\newcommand{\norm}[1]{\left\|#1\right\|}
\newcommand{\normtwo}[1]{\left\|#1\right\|_2}
\newcommand{\normF}[1]{\left\|#1\right\|_F}
\newcommand{\cO}{\mathcal{O}}
\newcommand{\J}{\left(\frac{1}{N}\one\one^\top\right)}
\newcommand{\Pmat}{I-\J}
\newcommand{\ProjPSD}{\Pi_{\succeq 0}}

\newcommand{\barx}{\bar x}
\newcommand{\barg}{\bar g}
\newcommand{\bars}{\bar s}

% ============================================================
% Title
% ============================================================
\title{\textbf{Distributed Gradient--Regularized Newton:\\
		Adaptive Consensus and \texorpdfstring{$\mathcal{O}(1/k^2)$}{O(1/k^2)} Global Convergence}}
\author{
	\textbf{Anonymous Authors} \\
	\textit{Institution} \\
	\texttt{email@domain.edu}
}
\date{\today}

%%%%%%%%%%%%%%%%%%%%%%%%%%%%%%%%%%%%%%%%%%%%%%%%%%%%%%%%%%%%
\begin{document}
	\maketitle
	
\begin{abstract}
	We study decentralized optimization of a convex finite-sum objective over a connected network
	of $N$ agents that communicate only with immediate neighbors.
	Second-order methods are appealing in communication-limited environments because they reduce
	the number of outer iterations, yet they are challenging to deploy in fully decentralized
	settings due to (i) the cost of maintaining curvature information across the network,
	(ii) sensitivity to stepsize selection, and (iii) network-induced disagreement and tracking errors.
	Motivated by stepsize-free regularized Newton schemes \citep{mishchenko2023rn},
	we propose \textsc{DisGrem}, a decentralized gradient--regularized Newton method that combines
	gradient tracking, Hessian tracking with PSD projection, and a two-stage mixing
	(pre-mixing before the local Newton solve, post-mixing after the update).
	A user-chosen scaling factor $M>0$ amplifies the vanishing regularization
	$\lambda_{i,k}=\sqrt{M\|\tilde g_{i,k}\|}$, enlarging the admissible heterogeneity regime
	while preserving $\lambda_{i,k}\to 0$ as the gradient vanishes.
	
	Our analysis introduces a \emph{reference-step} construction that resolves the
	non-commutativity between network averaging and matrix inversion, yielding an
	\emph{inexact regularized Newton} characterization of the average iterate.
	To quantify persistent curvature mismatch across agents, we introduce a bounded Hessian-heterogeneity
	constant over the relevant level set; it enters our network-dependent constants and is mild in near-i.i.d.\ regimes.
	Under convexity, Lipschitz Hessians, bounded level sets, and sufficiently accurate mixing,
	we establish $\cO(1/k^2)$ functional decay and $\cO(\varepsilon^{-1})$ iteration complexity
	to reach $\|\nabla f(\bar x_k)\|\le \varepsilon$ \emph{above a heterogeneity-dependent accuracy floor}
	(and recover the classical centralized guarantees in homogeneous/zero-bias regimes).
	Moreover, under a logarithmic accuracy schedule we obtain
	$\cO(\varepsilon^{-1}\log(\varepsilon^{-1}))$ total neighbor-communication rounds.
	
	We also develop practical variants: \textsc{CeDisGrem} compresses Hessian corrections
	to reduce communication by up to $50\%$, and \textsc{AdaDisGrem} adapts $M$ online.
	Extensive numerical experiments on 15 benchmark objectives, including ill-conditioned
	quadratics, log-sum-exp, real-data logistic regression, and non-convex landscapes,
	show that the \textsc{DisGrem} family converges to near machine precision $5$--$8\times$ faster
	than the best first-order baseline (EXTRA) and robustly handles weakly connected networks
	where gradient-tracking first-order methods fail entirely.
\end{abstract}
	
	%%%%%%%%%%%%%%%%%%%%%%%%%%%%%%%%%%%%%%%%%%%%%%%%%%%%%%%%%%%%
	\section{Introduction}
	%%%%%%%%%%%%%%%%%%%%%%%%%%%%%%%%%%%%%%%%%%%%%%%%%%%%%%%%%%%%
	Decentralized optimization addresses problems of the form
	\begin{equation}\label{eq:intro_prob}
		\min_{x\in\R^d} f(x)\;:=\;\frac{1}{N}\sum_{i=1}^N f_i(x),
	\end{equation}
	where each $f_i$ is stored privately at node $i$ of a connected network and nodes communicate only with neighbors.
	A principal bottleneck is \emph{communication}: algorithms must trade off local computation
	against costly neighbor exchanges.
	First-order decentralized methods may require many iterations on ill-conditioned problems, amplifying communication costs.
	Second-order information can reduce outer iterations, but decentralized second-order methods face major obstacles:
	curvature information is expensive to maintain and global convergence often requires careful stepsize control or line-search.
	
	A stepsize-free perspective is provided by the regularized Newton method of \citet{mishchenko2023rn},
	which uses a vanishing regularization scale $\lambda_k\asymp \sqrt{\|g_k\|}$ and achieves the optimal functional rate $\cO(1/k^2)$
	under convexity and Lipschitz Hessians without manual stepsize tuning.
	This motivates: \emph{can we obtain the same outer-iteration complexity in a decentralized network while controlling disagreement/tracking errors?}
	
\noindent\textbf{Contributions.}
We propose and analyze \textsc{DisGrem} (Distributed Gradient-Regularized Newton Method),
which combines gradient tracking, Hessian tracking with PSD projection, and two-stage
mixing (pre-mixing before the local Newton solve, post-mixing after the update).
Our specific contributions are:
\begin{itemize}[leftmargin=1.2em,itemsep=2pt]
	\item \textbf{Algorithm design.}
	We formulate a stepsize-free decentralized second-order method based on vanishing regularization
	$\lambda_{i,k}=\sqrt{M\|\tilde g_{i,k}\|}$.
	Two-stage mixing decouples the network averaging accuracy from the Newton solve,
	and PSD projection of tracked Hessians maintains positive semidefiniteness throughout.
	
	\item \textbf{Reference-step framework.}
	We introduce a reference-step construction to control the non-commutativity between network averaging
	and matrix inversion, yielding an \textbf{inexact regularized Newton} characterization for the average iterate.
	
	\item \textbf{Hessian heterogeneity analysis.}
	We identify the obstruction from persistent inter-agent Hessian mismatch and quantify it through
	a bounded heterogeneity constant $\sigma_H$ on the relevant level set; see~\eqref{eq:sigma_H_def}.
	
	\item \textbf{Convergence guarantees.}
	Under stated assumptions we prove $\Phi_k=\cO(1/k^2)$ functional decay and
	\textbf{$\cO(\varepsilon^{-1})$} iteration complexity to reach $\|\nabla f(\bar x_k)\|\le \varepsilon$
	for any $\varepsilon$ above the explicit floor in~\eqref{eq:eps_floor},
	and \textbf{$\cO(\varepsilon^{-1}\log(\varepsilon^{-1}))$} total mixing rounds under a logarithmic accuracy schedule.
	
	\item \textbf{Practical variants.}
	We develop \textbf{\textsc{CeDisGrem}}, which applies top-$r$ low-rank compression to Hessian correction messages
	and reduces per-iteration communication by 30--50\% with negligible iteration overhead;
	and \textbf{\textsc{AdaDisGrem}}, which adapts the scaling factor $M$ online to track local curvature changes.
	
	\item \textbf{Numerical evidence.}
	Experiments on 15 objectives across well-conditioned and ill-conditioned convex problems, real-data logistic regression,
	and non-convex landscapes show that the \textsc{DisGrem} family consistently converges to near machine precision while
	first-order baselines stagnate on weakly connected networks.
	It achieves $5$--$8\times$ fewer iterations than EXTRA and 100\% success rate in Monte-Carlo robustness trials
	where gradient-tracking baselines fail.
\end{itemize}
	
	%%%%%%%%%%%%%%%%%%%%%%%%%%%%%%%%%%%%%%%%%%%%%%%%%%%%%%%%%%%%
	\section{Related Work}
	%%%%%%%%%%%%%%%%%%%%%%%%%%%%%%%%%%%%%%%%%%%%%%%%%%%%%%%%%%%%
	
	\noindent\textbf{First-order decentralized methods.}
	Decentralized (sub)gradient descent dates to \citet{nedic2009subgradient}.
	The bias of constant-stepsize DGD is corrected by exact methods:
	EXTRA~\citep{shi2015extra} uses a two-step update with a carefully chosen pair of matrices,
	and exact diffusion~\citep{yuan2019exactdiffusion} reformulates the update in a primal-dual
	framework.
	Gradient tracking~\citep{nedic2017diging} achieves exact convergence
	under a single doubly stochastic matrix by maintaining a running sum of gradient increments.
	These methods achieve $\cO(1/k)$ or linear convergence for strongly convex objectives,
	but their iteration complexity depends on $(1-\rho)^{-1}$ or $(1-\rho)^{-2}$, which becomes
	prohibitive on weakly connected networks.
	Fast distributed gradient methods~\citep{jakovetic2014fast} improve over pure gradient descent
	but remain first-order.
	
	\noindent\textbf{Second-order and quasi-Newton decentralized methods.}
	Network Newton~\citep{mokhtari2015networknewton_arxiv} approximates the global Newton direction
	by a truncated power series of the Hessian, using multi-hop communication.
	Newton tracking~\citep{mokhtari2020newtontracking} integrates Newton-type directions into the
	tracking framework.
	Decentralized quasi-Newton methods, including DQM~\citep{eisen2017dqn}, replace the exact
	Hessian with L-BFGS-type approximations.
	All of these methods require explicit stepsize tuning and do not provide global convergence
	guarantees matching the centralized Newton rate.
	The challenge of Hessian heterogeneity has not been explicitly characterized in prior work.
	
	\noindent\textbf{Regularized Newton and cubic regularization.}
	The stepsize-free perspective adopted here is inspired by the centralized regularized Newton
	method of \citet{mishchenko2023rn}, which uses $\lambda_k\asymp\sqrt{\|\nabla f(x_k)\|}$
	and achieves the tight $\cO(1/k^2)$ functional rate under convexity and Lipschitz Hessians.
	This builds on cubic-regularization ideas of \citet{nesterov2006cubic,cartis2011arc},
	which also achieve $\cO(1/k^2)$ via an adaptive cubic model.
	Our paper is the first, to our knowledge, to show that the centralized $\cO(1/k^2)$
	rate can be preserved in a decentralized network via Hessian/gradient tracking and
	two-stage mixing, and the first to provide a transparent treatment of the Hessian
	heterogeneity obstruction.
	
	%%%%%%%%%%%%%%%%%%%%%%%%%%%%%%%%%%%%%%%%%%%%%%%%%%%%%%%%%%%%
	\section{Preliminaries and Assumptions}\label{sec:prelim}
	%%%%%%%%%%%%%%%%%%%%%%%%%%%%%%%%%%%%%%%%%%%%%%%%%%%%%%%%%%%%
	\subsection{Network model and notation}
	Communication uses a symmetric, nonnegative, doubly stochastic mixing matrix $W\in\R^{N\times N}$:
	$W_{ij}\ge0$, $W\one=\one$, and $\one^\top W=\one^\top$.
	Define the averaging projector $\J=\frac{1}{N}\one\one^\top$ and $\Pmat=I-\J$.
	Let $\rho := \normtwo{W-\J}\in[0,1)$, strictly less than $1$ for a connected undirected network.
	For stacked vectors $Z=[z_1;\dots;z_N]\in\R^{Nd}$ we use the Euclidean norm $\|Z\|$.
	We write $\bar z:=\frac1N\sum_{i=1}^N z_i$ for averages.
	
	\subsection{Standing assumptions}
	\begin{assumption}\label{ass:standing}
		\begin{enumerate}\itemsep2pt
			\item \textbf{Convexity and smoothness.}
			Each $f_i$ is convex and twice continuously differentiable.
			Its gradient is $L_1$-Lipschitz and its Hessian is $L_2$-Lipschitz:
			\[
			\norm{\nabla f_i(x)-\nabla f_i(y)}\le L_1\norm{x-y},
			\qquad
			\normtwo{\nabla^2 f_i(x)-\nabla^2 f_i(y)}\le L_2\norm{x-y}.
			\]
			\item \textbf{PSD Hessians.} $\nabla^2 f_i(x)\succeq0$ for all $x$ and all $i$.
			\item \textbf{Connected network.} $W$ is symmetric, doubly stochastic, and $\rho=\normtwo{W-\J}<1$.
			\item \textbf{Bounded level set (local iterates).} There exist $D>0$ and $x_\star\in\arg\min f$ such that for all $k\ge 0$ and all $i\in\{1,\dots,N\}$,
			\begin{equation}\label{eq:bounded_levelset}
				\norm{x_{i,k}-x_\star}\le D.
			\end{equation}
			Consequently, $\norm{\bar x_k-x_\star}\le D$ and $\norm{\tilde x_{i,k}-x_\star}\le D$ for all $i,k$ because $\bar x_k$ and $\tilde x_{i,k}$ are convex combinations of the $x_{j,k}$.
			\item \textbf{Bounded Hessians on the relevant set.}
			There exists $M_H\ge0$ such that $\normtwo{\nabla^2 f(x)}\le M_H$ for all $x$ with $\norm{x-x_\star}\le D$.
			\item \textbf{Bounded Hessian heterogeneity (level-set constant).}
			Define
			\begin{equation}\label{eq:sigma_H_def}
				\sigma_H := \sup\Big\{\, \frac{1}{\sqrt N}\normF{((\Pmat)\otimes I_{d^2})\,\nabla^2 F(\one\otimes x)}\;:\; \norm{x-x_\star}\le D\,\Big\}\in[0,\infty).
			\end{equation}
			This constant measures the RMS mismatch of local Hessians relative to the average Hessian when evaluated at the same point.
			\item \textbf{Uniform bound on averaged Hessian-mismatch term.} Define the per-node mismatch
			$\hat e_{i,k}:=\normF{\tilde H_{i,k}-\nabla^2 f_i(\tilde x_{i,k})}$ and its average $\bar e_k:=\frac1N\sum_{i=1}^N\hat e_{i,k}$.
			There exists a finite constant $\bar e_{\max}\ge 0$ such that
			\begin{equation}\label{eq:ebar_max}
				\bar e_k\le \bar e_{\max}\qquad\text{for all }k\ge 0.
			\end{equation}
			This assumption can be enforced in practice by clipping $\hat e_{i,k}$ or by spectrally truncating tracked Hessians/corrections.
			\item \textbf{Regularization scaling.} The user-chosen scaling factor satisfies $M\ge L_2$.
		\end{enumerate}
	\end{assumption}
	
	\begin{remark}[Interpretation of Assumption~\ref{ass:standing}(6)]
		Assumption~\ref{ass:standing}(6) is a \emph{boundedness} condition on curvature mismatch over the relevant level set.
		It is automatically satisfied whenever $\nabla^2 f_i$ are bounded on $\{x:\ \norm{x-x_\star}\le D\}$, and it becomes small in near-i.i.d.\ regimes.
		Unlike a vanishing condition, it does \emph{not} force local Hessians to coincide at stationary points; rather, it quantifies the irreducible mismatch that must be controlled by mixing and regularization.
	\end{remark}
	
	% ============================================================
	% Derived heterogeneity constant (finite on the bounded level set)
	% ============================================================
	Define the (possibly zero) gradient-heterogeneity constant
	\begin{equation}\label{eq:sigma_g_def}
		\sigma_g := \sup\Big\{\, \norm{((\Pmat)\otimes I_d)\,\nabla F(\one\otimes x)}\;:\; \norm{x-x_\star}\le D\,\Big\}\in[0,\infty),
	\end{equation}
	which is finite under Assumption~\ref{ass:standing}(1) and (4) because $\nabla F(\one\otimes x)$ is continuous on a compact set.
	In interpolation/shared-minimizer regimes one typically has $\sigma_g=0$, while for general heterogeneous data this constant may be positive and induces an unavoidable accuracy floor.
	
	
	%%%%%%%%%%%%%%%%%%%%%%%%%%%%%%%%%%%%%%%%%%%%%%%%%%%%%%%%%%%%
	\section{Algorithm}\label{sec:alg}
	%%%%%%%%%%%%%%%%%%%%%%%%%%%%%%%%%%%%%%%%%%%%%%%%%%%%%%%%%%%%
	
	\subsection{Design motivation}\label{ssec:design}
	
	Three interacting challenges must be resolved to lift regularized Newton to the decentralized setting.
	
	\textbf{(i) Averaging vs.\ inversion.}
	In the centralized regularized Newton method of \citet{mishchenko2023rn} the Newton system involves
	$(\nabla^2 f(\bar x_k)+\lambda_k I)^{-1}$, which is an invertible function of the \emph{average} gradient.
	In a network each agent solves its own local system, so the true average step
	$\bar s_k = \frac1N\sum_i s_{i,k} \ne (\frac1N\sum_i A_{i,k})^{-1}(\frac1N\sum_i g_{i,k})$
	in general: averaging and inversion do not commute.
	Pre-mixing reduces disagreement in the inputs $(x_{i,k}, g_{i,k}, H_{i,k})$ before the solve,
	making individual systems closer to the average system and bounding the resulting
	\emph{step dispersion} (Section~\ref{sec:theory}, Lemma~\ref{lem:step-disp}).
	
	\textbf{(ii) Hessian maintenance and positivity.}
	Tracker dynamics propagate Hessian increments $\nabla^2 f_j(x_{j,k+1})-\nabla^2 f_j( x_{j,k})$
	through the network.
	Floating-point errors and network delays may produce non-symmetric or indefinite tracked Hessians,
	which would make the regularized Newton system ill-conditioned.
	We therefore symmetrize and project $\widehat H_{i,k+1}$ onto the PSD cone in Frobenius norm
	before storing $H_{i,k+1}$.
	This projection is non-expansive (Remark~\ref{rem:psd-nonexp}) and introduces an error
	$\hat e_{i,k}:=\normF{\tilde H_{i,k}-\nabla^2 f_i(\tilde x_{i,k})}$ that is absorbed into
	the regularization via the augmented shift $\lambda_{i,k}+\hat e_{i,k}$,
	guaranteeing a well-posed local solve for any tracked Hessian.
	
	\textbf{(iii) Stepsize-free regularization.}
	Setting $\lambda_{i,k}=\sqrt{M\|\tilde g_{i,k}\|}$ mirrors the centralized choice
	$\lambda_k\asymp\sqrt{\|\nabla f(\bar x_k)\|}$, which achieves $\cO(1/k^2)$ without a line search
	or manual stepsize.
	The user constant $M\ge L_2$ scales the regularizer uniformly: larger $M$ produces larger shifts,
	which damp Newton steps, tolerate wider heterogeneity, and improve stability on
	ill-conditioned or non-convex problems (Section~\ref{ssec:param});
	once the iterates enter a region where $\|\nabla f\|\to 0$, the regularizer vanishes automatically.
	
	\subsection{\textsc{DisGrem}: full algorithm}\label{ssec:disgrem}
	
	Each node $i$ stores a primal variable $x_{i,k}\in\R^d$, a gradient tracker $g_{i,k}\in\R^d$,
	and a Hessian tracker $H_{i,k}\in\R^{d\times d}$.
	$\ProjPSD(\cdot)$ denotes the Euclidean (Frobenius) projection onto the PSD cone.
	
	\begin{algorithm}[H]
		\caption{Decentralized Gradient--Regularized Newton (\textsc{DisGrem})}
		\label{alg:disgrem}
		\begin{algorithmic}[1]
			\State \textbf{Input:} $\{x_{i,0}\}_{i=1}^N$, mixing matrix $W$, Lipschitz constants $L_1,L_2>0$, scaling factor $M\ge L_2$.
			\State \textbf{Initialize:} $g_{i,0}\gets \nabla f_i(x_{i,0})$, $H_{i,0}\gets \nabla^2 f_i(x_{i,0})$.
			\For{$k=0,1,2,\dots$}
			\State \textbf{(A) Pre-mixing:} choose $\tau_k\ge 1$ mixing rounds and set
			\[
			\tilde x_{i,k}\gets \sum_{j=1}^N [W^{\tau_k}]_{ij}\,x_{j,k},\quad
			\tilde g_{i,k}\gets \sum_{j=1}^N [W^{\tau_k}]_{ij}\,g_{j,k},\quad
			\tilde H_{i,k}\gets \sum_{j=1}^N [W^{\tau_k}]_{ij}\,H_{j,k}.
			\]
			\State \textbf{(B) Local computation (parallel for each $i$):}
			\[
			\lambda_{i,k} \gets \sqrt{M\,\norm{\tilde g_{i,k}}},\quad
			\hat e_{i,k} \gets \normF{\tilde H_{i,k}-\nabla^2 f_i(\tilde x_{i,k})}.
			\]
			\Statex\hspace{1.6em}\textbf{Solve} $(\tilde H_{i,k}+(\lambda_{i,k}+\hat e_{i,k})I)\,s_{i,k}=-\tilde g_{i,k}$
			\Statex\hspace{1.6em}(set $s_{i,k}=0$ if $\tilde g_{i,k}=0$); set $y_{i,k+1}\gets \tilde x_{i,k}+s_{i,k}$.
			\State \textbf{(C) Post-mixing:} choose $t_k\ge 1$ mixing rounds and set
			\[
			x_{i,k+1}\gets \sum_{j=1}^N [W^{t_k}]_{ij}\,y_{j,k+1}.
			\]
	\State \textbf{(D) Tracker updates:}
	\[
	g_{i,k+1}\gets \sum_{j=1}^N [W^{t_k}]_{ij}\!\Big(\tilde g_{j,k}+\nabla f_j(x_{j,k+1})-\nabla f_j(x_{j,k})\Big),
	\]
	\[
	\widehat H_{i,k+1}\gets \sum_{j=1}^N [W^{t_k}]_{ij}\!\Big(\tilde H_{j,k}+\nabla^2 f_j(x_{j,k+1})-\nabla^2 f_j(x_{j,k})\Big),
	\]
	\[
	H_{i,k+1}\gets \ProjPSD\!\left(\tfrac12\bigl(\widehat H_{i,k+1}+\widehat H_{i,k+1}^\top\bigr)\right).
	\]
			\EndFor
		\end{algorithmic}
	\end{algorithm}
	
\subsection{Communication cost per iteration}\label{ssec:comm}

Each outer iteration of \textsc{DisGrem} involves $(\tau_k + t_k)$ rounds of neighbor communication.
In each round every agent $i$ sends messages to and receives messages from its $|\mathcal N_i|$ neighbors.
In \textsc{DisGrem}, the natural payload exchanged per neighbor per round consists of:
\begin{itemize}[leftmargin=1.2em,itemsep=1pt]
	\item two $d$-dimensional vectors (e.g., mixing of $(x,g)$ in pre-mixing, or mixing of $(y,g)$ in post-mixing): $2d$ floats,
	\item one $d\times d$ symmetric matrix (Hessian tracker or Hessian increment): $\tfrac{d(d+1)}{2}$ floats.
\end{itemize}
Equivalently, the two vectors can be concatenated into a single $2d$-vector; this does not change the byte complexity.

Let $n_v$ denote the number of $d$-vectors transmitted per round (for \textsc{DisGrem}, $n_v=2$).
With double precision (8 bytes per float), the \emph{sent} message volume per agent per outer iteration is
\begin{equation}\label{eq:comm_cost_iter}
	C_{\mathrm{iter}}^{\mathrm{send}}(k)
	= (\tau_k+t_k)\cdot |\mathcal N_i|\cdot \Bigl(n_v d + \tfrac{d(d+1)}{2}\Bigr)\cdot 8 \ \ \text{bytes}.
\end{equation}
If one counts both sending and receiving volume, then
$C_{\mathrm{iter}}^{\mathrm{send+recv}}(k)=2\,C_{\mathrm{iter}}^{\mathrm{send}}(k)$.

For the default setting $N=10$, $d=10$, and the random geometric graph with $r_0=0.5$ giving average degree
$|\mathcal N_i|\approx 3.5$, taking $\tau_k=t_k=3$ and $n_v=2$ yields
\[
C_{\mathrm{iter}}^{\mathrm{send}}\approx 6\times 3.5\times\Bigl(20+55\Bigr)\times 8
\approx 12.6\,\text{kB}
\quad \text{per agent per outer iteration}.
\]
This back-of-the-envelope estimate is consistent with the empirical \texttt{commCost} magnitudes reported in
Table~\ref{tab:main} (up to graph-degree variability and implementation details such as message aggregation and bookkeeping).
	
	\subsection{Practical variants}\label{ssec:variants}
	
\noindent\textbf{\textsc{CeDisGrem}: compressed Hessian communication.}
The dominant communication cost of \textsc{DisGrem} is the $\frac{d(d+1)}{2}$-dimensional
Hessian increment $\nabla^2 f_j(x_{j,k+1})-\nabla^2 f_j(x_{j,k})$.
\textsc{CeDisGrem} replaces each symmetric Hessian increment
$\Delta H_{j,k}:=\nabla^2 f_j(x_{j,k+1})-\nabla^2 f_j(x_{j,k})$
by a rank-$r$ symmetric truncation based on the eigen-decomposition.
Let $\Delta H_{j,k}=V\operatorname{diag}(\mu_1,\dots,\mu_d)V^\top$ with
$|\mu_1|\ge \cdots \ge |\mu_d|$ and $V=[v_1,\dots,v_d]$ orthogonal.
Define the truncated approximation
\[
\hat\Delta H_{j,k}:=\sum_{\ell=1}^r \mu_\ell v_\ell v_\ell^\top
= V_r\Lambda_r V_r^\top,
\]
where $V_r=[v_1,\dots,v_r]\in\R^{d\times r}$ and $\Lambda_r=\operatorname{diag}(\mu_1,\dots,\mu_r)\in\R^{r\times r}$.
This $\hat\Delta H_{j,k}$ is the best rank-$r$ symmetric approximation of $\Delta H_{j,k}$ in Frobenius norm
(equivalently, the best rank-$r$ approximation for a symmetric matrix is obtained by keeping the $r$
eigenpairs with largest absolute eigenvalues).
Thus, transmitting $(V_r,\Lambda_r)$ costs $r(d+1)$ floats instead of $\frac{d(d+1)}{2}$ floats per neighbor per round.

The truncation error satisfies
\[
\normF{\Delta H_{j,k}-\hat\Delta H_{j,k}}
=\left(\sum_{\ell=r+1}^d \mu_\ell^2\right)^{1/2}
\le |\mu_{r+1}(\Delta H_{j,k})|\sqrt{d-r},
\]
where $\mu_{r+1}(\Delta H_{j,k})$ denotes the $(r+1)$-th largest absolute eigenvalue.
This error is absorbed into $\hat e_{i,k}$ and thus into the regularization automatically.
Empirically, $r=\lceil d/5\rceil$ reduces communication by 30--50\% with negligible iteration overhead
(Table~\ref{tab:main}).
	
	\noindent\textbf{\textsc{AdaDisGrem}: adaptive scaling factor.}
	The default choice $M=\mathrm{const}$ requires prior knowledge of $L_2$.
	\textsc{AdaDisGrem} maintains a running estimate $\hat M_k$ of the local Hessian Lipschitz
	constant using the secant condition:
	\[
	\hat M_k \gets \max\!\Bigl(\hat M_{k-1},\;
	\frac{\normtwo{\nabla^2 f_i(x_{i,k})-\nabla^2 f_i(x_{i,k-1})}}{\norm{x_{i,k}-x_{i,k-1}}}\Bigr),
	\]
	and sets $\lambda_{i,k}=\sqrt{\hat M_k\|\tilde g_{i,k}\|}$.
	This heuristic removes the need to specify $L_2$ in practice; a full theoretical treatment of the adaptive rule is left for future work.
	
	\subsection{Practical guidelines for parameter selection}\label{ssec:practical}
	
	\begin{itemize}[leftmargin=1.2em,itemsep=2pt]
		\item \textbf{Scaling factor $M$.}
		Theorem~\ref{thm:main} requires $M\ge L_2$.
		In practice, setting $M = c \cdot \hat L$, where $\hat L$ is a rough estimate of the
		gradient Lipschitz constant and $c\in[0.5, 2.0]$, works well across all tested problems.
		Our parameter study (Section~\ref{ssec:param}) shows convergence is largely insensitive
		to the exact value of $M$ within this range.
		When $L_2$ is unknown, use \textsc{AdaDisGrem}.
		\item \textbf{Consensus rounds $N_c$.}
		The theory requires $\tau_k, t_k = \cO(\log 1/\|\tilde{\bar g}_k\|)$, which grows
		logarithmically.  In practice, a fixed $N_c = 3$--$5$ rounds per stage is sufficient
		for most networks with $\rho \le 0.95$.
		For strongly sparse networks ($\rho > 0.97$), increasing $N_c$ to $8$--$10$
		significantly improves stability.
		\item \textbf{Hessian compression rank $r$ (\textsc{CeDisGrem}).}
		Set $r = \lceil 0.2d \rceil$ as a default; increase $r$ if convergence slows.
		For diagonal-dominant Hessians (e.g.\ $\ell_2$-regularized losses), even $r=1$
		is often sufficient.
	\end{itemize}
	
	%%%%%%%%%%%%%%%%%%%%%%%%%%%%%%%%%%%%%%%%%%%%%%%%%%%%%%%%%%%%
	\section{Theory and Convergence Analysis}\label{sec:theory}
	%%%%%%%%%%%%%%%%%%%%%%%%%%%%%%%%%%%%%%%%%%%%%%%%%%%%%%%%%%%%
	Throughout, Assumption~\ref{ass:standing} holds.
	All proofs are provided in Appendix~\ref{app:proofs}.
	
	\subsection{Basic inequalities}
	\begin{lemma}[Third-order Taylor bound]\label{lem:taylor3}
		Let $h:\R^d\to\R$ have $L_2$-Lipschitz Hessian. Then for all $x,s$,
		\[
		h(x+s)\le h(x)+\ip{\nabla h(x)}{s}+\frac12 s^\top\nabla^2 h(x)s+\frac{L_2}{6}\norm{s}^3.
		\]
	\end{lemma}
	
	\begin{lemma}[Consensus contraction]\label{lem:consensus}
		For any $z\in\R^N$ and integer $t\ge 1$,
		\[
		\normtwo{(W^t-\J)z}\le \rho^t\normtwo{z}.
		\]
		Consequently, for any stacked $Z\in\R^{Nd}$,
		\[
		\norm{((W^t-\J)\otimes I_d)Z}\le \rho^t\norm{Z}.
		\]
	\end{lemma}
	
	\begin{lemma}[Smooth convex gap--gradient inequality]\label{lem:gap-grad}
		Let $h$ be convex with $L$-Lipschitz gradient. Then for any minimizer $x_\star$,
		\[
		\norm{\nabla h(x)}^2\le 2L\big(h(x)-h(x_\star)\big).
		\]
	\end{lemma}
	
	\subsection{Frobenius stabilizer: a local descent inequality}
	\begin{lemma}[Frobenius domination]\label{lem:frob}
		Let $E$ be symmetric and $\hat e=\normF{E}$. Then $\hat e\ge \normtwo{E}$ and $\hat e I\pm E\succeq 0$.
		In particular, for any $s$, $-\hat e\norm{s}^2-s^\top E s\le 0$.
	\end{lemma}
	
	\begin{lemma}[Scaled regularization dominates step size]\label{lem:M-vs-s}
		Assume $\tilde H_{i,k}\succeq 0$ and $M\ge L_2$. Then
		\[
		M\norm{s_{i,k}}\le \lambda_{i,k}
		\qquad\text{and hence}\qquad
		L_2\norm{s_{i,k}}\le \lambda_{i,k}.
		\]
	\end{lemma}
	
	\begin{lemma}[Local descent with tracking]\label{lem:local-descent}
		Fix $i,k$. Define $\tilde\delta_{i,k}:=\tilde g_{i,k}-\nabla f_i(\tilde x_{i,k})$ and
		$\tilde e_{i,k}:=\tilde H_{i,k}-\nabla^2 f_i(\tilde x_{i,k})$.
		Then
		\begin{equation}\label{eq:local-descent}
			f_i(\tilde x_{i,k}+s_{i,k})
			\le f_i(\tilde x_{i,k})
			-\lambda_{i,k}\norm{s_{i,k}}^2
			-\ip{\tilde\delta_{i,k}}{s_{i,k}}
			+\frac{L_2}{6}\norm{s_{i,k}}^3.
		\end{equation}
	\end{lemma}
	
	\begin{lemma}[Clean local descent]\label{lem:local-clean}
		Under Lemma~\ref{lem:local-descent} and Lemma~\ref{lem:M-vs-s},
		\[
		f_i(\tilde x_{i,k}+s_{i,k})
		\le f_i(\tilde x_{i,k})
		-\frac{5}{6}\lambda_{i,k}\norm{s_{i,k}}^2
		-\ip{\tilde\delta_{i,k}}{s_{i,k}}.
		\]
	\end{lemma}
	
	\subsection{Average dynamics and reference step}
	Define pre-mixed averages:
	\[
	\tilde{\bar x}_k:=\frac1N\sum_{i=1}^N\tilde x_{i,k},\quad
	\tilde{\bar g}_k:=\frac1N\sum_{i=1}^N\tilde g_{i,k},\quad
	\tilde{\bar H}_k:=\frac1N\sum_{i=1}^N\tilde H_{i,k},
	\]
	and the average step $\bar s_k:=\frac1N\sum_{i=1}^N s_{i,k}$.
	
	\begin{lemma}[Average preservation]\label{lem:avg-pres}
		For all $k$, $\tilde{\bar x}_k=\bar x_k$ and $\bar x_{k+1}=\bar x_k+\bar s_k$.
	\end{lemma}
	
\begin{lemma}[Average tracking identity]\label{lem:avg-track}
	Assume Algorithm~\ref{alg:disgrem} uses the tracker update in step \textbf{(D)} and is initialized by
	$g_{i,0}=\nabla f_i(x_{i,0})$.
	Then for all $k\ge 0$,
	\[
	\bar g_k:=\frac1N\sum_{i=1}^N g_{i,k}=\frac1N\sum_{i=1}^N \nabla f_i(x_{i,k}),
	\qquad
	\tilde{\bar g}_k=\bar g_k.
	\]
\end{lemma}
\begin{proof}
	Because $W^{t_k}$ is doubly stochastic, averaging step \textbf{(D)} gives
	\[
	\bar g_{k+1}=\tilde{\bar g}_k+\frac1N\sum_{i=1}^N\big(\nabla f_i(x_{i,k+1})-\nabla f_i(x_{i,k})\big).
	\]
	Also $\tilde{\bar g}_k=\bar g_k$ since pre-mixing preserves averages.
	By induction with $\bar g_0=\frac1N\sum_i\nabla f_i(x_{i,0})$, the claim follows.
\end{proof}
	
	Define
	\[
	\tilde\lambda_{i,k}:=\lambda_{i,k}+\hat e_{i,k},\qquad
	\tilde{\bar\lambda}_k:=\frac1N\sum_{i=1}^N\tilde\lambda_{i,k},\qquad
	A_{i,k}:=\tilde H_{i,k}+\tilde\lambda_{i,k}I.
	\]
	Then $s_{i,k}=-A_{i,k}^{-1}\tilde g_{i,k}$.
	Define the reference system
	\[
	A_k^{\mathrm{ref}}:=\tilde{\bar H}_k+\tilde{\bar\lambda}_k I,\qquad
	s_k^{\mathrm{ref}}:=-(A_k^{\mathrm{ref}})^{-1}\tilde{\bar g}_k.
	\]
	
	\begin{lemma}[Average regularization regime]\label{lem:avg-regime}
		For all $k$,
		\[
		L_2\norm{\bar s_k}\le \frac1N\sum_{i=1}^N \lambda_{i,k}\le \tilde{\bar\lambda}_k.
		\]
	\end{lemma}
	
	\subsection{Bridge bounds and step dispersion}
	Define the pre-mixed disagreement energy
	\[
	\tilde E_k:=\frac{1}{N}\sum_{i=1}^N \norm{\tilde x_{i,k}-\bar x_k}^2.
	\]
	Define the post-mixing primal dispersion
	\[
	D_X(k):=\sqrt{\frac1N\sum_{i=1}^N \|x_{i,k}-\bar x_k\|^2}.
	\]
	Define the pre-mixed gradient tracking residuals
	\[
	\tilde\delta_{i,k}:=\tilde g_{i,k}-\nabla f_i(\tilde x_{i,k}),\qquad
	\tilde\Delta_k:=[\tilde\delta_{1,k};\dots;\tilde\delta_{N,k}]\in\R^{Nd}.
	\]
	
	
\begin{lemma}[A computable upper bound for $\norm{\tilde\Delta_k}$]\label{lem:tildeDelta-upper}
	For every $k$,
	\[
	\frac{\norm{\tilde\Delta_k}}{\sqrt N}
	\le D(\tilde G_k) + L_1\sqrt{\tilde E_k} + L_1 D_X(k) + \frac{\sigma_g}{\sqrt N},
	\qquad\text{where }\;D(\tilde G_k):=\sqrt{\frac1N\sum_{i=1}^N\norm{\tilde g_{i,k}-\tilde{\bar g}_k}^2}.
	\]
	In particular, $\Delta^g_k\le D(\tilde G_k)$ and hence
	$\frac{\norm{\tilde\Delta_k}}{\sqrt N}\le \Delta^g_k + L_1\sqrt{\tilde E_k} + L_1 D_X(k) + \frac{\sigma_g}{\sqrt N}$.
\end{lemma}
	
	\begin{lemma}[Gradient bridge]\label{lem:bridge-grad}
		For every $k$,
		\[
		\norm{\nabla f(\bar x_k)-\tilde{\bar g}_k}
		\le \frac{\norm{\tilde\Delta_k}}{\sqrt N}+L_1\sqrt{\tilde E_k}.
		\]
	\end{lemma}
	
	\begin{lemma}[Hessian bridge]\label{lem:bridge-hess}
		For every $k$,
		\[
		\normtwo{\nabla^2 f(\bar x_k)-\tilde{\bar H}_k}
		\le L_2\sqrt{\tilde E_k}+\bar e_k,\qquad
		\bar e_k:=\frac1N\sum_{i=1}^N \hat e_{i,k}.
		\]
	\end{lemma}
	
	\begin{remark}[PSD projection is non-expansive]\label{rem:psd-nonexp}
		$\ProjPSD(\cdot)$ is the Euclidean (Frobenius) projection onto a closed convex cone, hence
		$\normF{\ProjPSD(A)-\ProjPSD(B)}\le \normF{A-B}$ for all $A,B$.
	\end{remark}
	
	\begin{lemma}[PSD projection does not increase dispersion]\label{lem:psd-dispersion}
		Let $T(A):=\ProjPSD\bigl(\tfrac12(A+A^\top)\bigr)$. For any collection $\{A_i\}_{i=1}^N\subset\R^{d\times d}$,
		if we set $B_i:=T(A_i)$, then
		\[
		\frac1N\sum_{i=1}^N\normF{B_i-\bar B}^2
		\le \frac1N\sum_{i=1}^N\normF{A_i-\bar A}^2,
		\qquad\bar A:=\frac1N\sum_i A_i,\quad \bar B:=\frac1N\sum_i B_i.
		\]
		Equivalently, the Frobenius dispersion $D_H(B):=\sqrt{\frac1N\sum_i\normF{B_i-\bar B}^2}$ satisfies $D_H(B)\le D_H(A)$.
	\end{lemma}
	
	
	Define dispersion measures
	\[
	\Delta^g_k:=\frac1N\sum_{i=1}^N \norm{\tilde g_{i,k}-\tilde{\bar g}_k},\quad
	\Delta^H_k:=\frac1N\sum_{i=1}^N \normtwo{\tilde H_{i,k}-\tilde{\bar H}_k},\quad
	\Delta^\lambda_k:=\frac1N\sum_{i=1}^N \big|\tilde\lambda_{i,k}-\tilde{\bar\lambda}_k\big|,
	\]
	and $\underline\lambda_k:=\min_{1\le i\le N}\tilde\lambda_{i,k}$.
	
	\begin{lemma}[Step dispersion via resolvent identity]\label{lem:step-disp}
		For every $k$,
		\[
		\norm{\bar s_k-s^{\mathrm{ref}}_k}
		\le \frac{1}{\underline\lambda_k}\Delta^g_k
		+\frac{\norm{\tilde{\bar g}_k}}{\underline\lambda_k\,\tilde{\bar\lambda}_k}\big(\Delta^H_k+\Delta^\lambda_k\big).
		\]
	\end{lemma}
\begin{lemma}[Control of the regularization dispersion]\label{lem:delta_lambda_control}
	Recall $\tilde\lambda_{i,k}=\lambda_{i,k}+\hat e_{i,k}$ and
	$\tilde{\bar\lambda}_k=\frac1N\sum_{i=1}^N \tilde\lambda_{i,k}$, where
	$\lambda_{i,k}=\sqrt{M\|\tilde g_{i,k}\|}$ and $\bar e_k:=\frac1N\sum_{i=1}^N \hat e_{i,k}$.
Recall the dispersions $\Delta_k^\lambda$ and $\Delta_k^g$ defined above.
	Assume the relative dispersion condition
	\[
	\max_{1\le i\le N}\|\tilde g_{i,k}-\tilde{\bar g}_k\|\le \alpha_d\|\tilde{\bar g}_k\|
	\quad\text{for some }\alpha_d\in(0,1).
	\]
	Then for any $k$ with $\tilde{\bar g}_k\neq 0$,
	\begin{equation}\label{eq:delta_lambda_bound}
		\Delta_k^\lambda
		\;\le\;
		\frac{\sqrt{M}}{\sqrt{1-\alpha_d}}\cdot
		\frac{\Delta_k^g}{\sqrt{\|\tilde{\bar g}_k\|}}
		\;+\; 2\bar e_k.
	\end{equation}
\end{lemma}

\begin{proof}
	Write $\bar\lambda_k:=\frac1N\sum_{i=1}^N \lambda_{i,k}$ so that
	$\tilde{\bar\lambda}_k=\bar\lambda_k+\bar e_k$.
	By the triangle inequality,
	\[
	\big|\tilde\lambda_{i,k}-\tilde{\bar\lambda}_k\big|
	=
	\big|(\lambda_{i,k}-\bar\lambda_k)+(\hat e_{i,k}-\bar e_k)\big|
	\le
	|\lambda_{i,k}-\bar\lambda_k|+|\hat e_{i,k}-\bar e_k|.
	\]
	Averaging over $i$ yields
	\begin{equation}\label{eq:dl_split}
		\Delta_k^\lambda \le
		\underbrace{\frac1N\sum_{i=1}^N |\lambda_{i,k}-\bar\lambda_k|}_{=:T_1}
		+
		\underbrace{\frac1N\sum_{i=1}^N |\hat e_{i,k}-\bar e_k|}_{=:T_2}.
	\end{equation}
	
	\noindent\emph{Step 1: bound $T_1$.}
	For any scalars $\{u_i\}_{i=1}^N$ with $\bar u=\frac1N\sum_i u_i$,
	\[
	\frac1N\sum_{i=1}^N |u_i-\bar u|
	\le \frac{1}{N^2}\sum_{i=1}^N\sum_{j=1}^N |u_i-u_j|,
	\]
	because $|u_i-\bar u|=\big|\frac1N\sum_j(u_i-u_j)\big|\le \frac1N\sum_j|u_i-u_j|$.
	Apply this with $u_i=\lambda_{i,k}$ to obtain
	\[
	T_1\le \frac{1}{N^2}\sum_{i=1}^N\sum_{j=1}^N |\lambda_{i,k}-\lambda_{j,k}|.
	\]
	Let $a_i:=\|\tilde g_{i,k}\|$. Then
	\[
	|\lambda_{i,k}-\lambda_{j,k}|
	=\sqrt{M}\,|\sqrt{a_i}-\sqrt{a_j}|
	=\sqrt{M}\,\frac{|a_i-a_j|}{\sqrt{a_i}+\sqrt{a_j}}.
	\]
	Moreover $|a_i-a_j|=\big|\|\tilde g_{i,k}\|-\|\tilde g_{j,k}\|\big|\le \|\tilde g_{i,k}-\tilde g_{j,k}\|$.
	Using $\|\tilde g_{i,k}-\tilde g_{j,k}\|\le \|\tilde g_{i,k}-\tilde{\bar g}_k\|+\|\tilde g_{j,k}-\tilde{\bar g}_k\|$ and averaging over $(i,j)$ gives
	\[
	\frac{1}{N^2}\sum_{i,j}\|\tilde g_{i,k}-\tilde g_{j,k}\|
	\le \frac{1}{N^2}\sum_{i,j}\Big(\|\tilde g_{i,k}-\tilde{\bar g}_k\|+\|\tilde g_{j,k}-\tilde{\bar g}_k\|\Big)
	=2\Delta_k^g.
	\]
	Under the relative dispersion assumption,
	$a_i\ge \|\tilde{\bar g}_k\|-\|\tilde g_{i,k}-\tilde{\bar g}_k\|\ge (1-\alpha_d)\|\tilde{\bar g}_k\|$,
	hence $\sqrt{a_i}+\sqrt{a_j}\ge 2\sqrt{(1-\alpha_d)\|\tilde{\bar g}_k\|}$ for all $(i,j)$.
	Therefore,
	\[
	T_1\le
	\sqrt{M}\cdot
	\frac{2\Delta_k^g}{2\sqrt{(1-\alpha_d)\|\tilde{\bar g}_k\|}}
	=
	\frac{\sqrt{M}}{\sqrt{1-\alpha_d}}\cdot
	\frac{\Delta_k^g}{\sqrt{\|\tilde{\bar g}_k\|}}.
	\]
	
	\noindent\emph{Step 2: bound $T_2$.}
	Since $\hat e_{i,k}\ge 0$ for all $i$,
	\[
	\sum_{i=1}^N |\hat e_{i,k}-\bar e_k|
	=
	2\sum_{i:\ \hat e_{i,k}>\bar e_k} (\hat e_{i,k}-\bar e_k)
	\le
	2\sum_{i:\ \hat e_{i,k}>\bar e_k}\hat e_{i,k}
	\le 2\sum_{i=1}^N \hat e_{i,k}=2N\bar e_k,
	\]
	so $T_2\le 2\bar e_k$.
	Combine with \eqref{eq:dl_split}.
\end{proof}
	
	\subsection{Inexact regularized Newton condition for the average step}
	Set $\eta:=1/12$ and define
	\[
	r_k := (\nabla^2 f(\bar x_k)+\tilde{\bar\lambda}_k I)\bar s_k+\nabla f(\bar x_k).
	\]
	
	\begin{proposition}[Inexact RN condition for the average step]\label{prop:inexact-rn}
		Set $\eta:=1/12$.
		Assume the following hold at iteration $k$:
		\begin{enumerate}[leftmargin=1.2em,itemsep=2pt]
			\item \textbf{Dispersion control:}
			\[
			\norm{\bar s_k-s^{\mathrm{ref}}_k}\le \frac{\eta}{8}\cdot \frac{\tilde{\bar\lambda}_k}{M_H+\tilde{\bar\lambda}_k}\,\norm{\bar s_k}.
			\]
			\item \textbf{Gradient bridge accuracy:}
			\[
			\frac{\norm{\tilde\Delta_k}}{\sqrt N}+L_1\sqrt{\tilde E_k}\le \frac{\eta}{8}\,\tilde{\bar\lambda}_k\norm{\bar s_k}.
			\]
			\item \textbf{Hessian bridge accuracy:}
			\[
			L_2\sqrt{\tilde E_k}+\bar e_k\le \frac{\eta}{8}\,\tilde{\bar\lambda}_k.
			\]
			\item \textbf{Hessian bound:} $\normtwo{\nabla^2 f(\bar x_k)}\le M_H$ (automatic by Assumption~\ref{ass:standing}(5)).
		\end{enumerate}
		Then
		\[
		\norm{r_k}\le \eta\,\tilde{\bar\lambda}_k\norm{\bar s_k},
		\qquad\text{and}\qquad
		L_2\norm{\bar s_k}\le \tilde{\bar\lambda}_k.
		\]
	\end{proposition}
	
	\subsection{Regularization comparability and heterogeneity absorption}
	Let $g_k:=\nabla f(\bar x_k)$ and $\lambda_k:=\tilde{\bar\lambda}_k$.
	
	\begin{lemma}[Comparability of $\lambda_k$ to $\sqrt{M\|g_k\|}$]\label{lem:lambda-comparable}
		Assume there exist $\alpha_b\in(0,1)$ and $\alpha_d\in(0,1)$ such that
		\begin{align}
			\label{eq:rel-bridge}
			\norm{g_k-\tilde{\bar g}_k} &\le \alpha_b \norm{\tilde{\bar g}_k},\\
			\label{eq:rel-disp}
			\max_{1\le i\le N}\norm{\tilde g_{i,k}-\tilde{\bar g}_k} &\le \alpha_d \norm{\tilde{\bar g}_k}.
		\end{align}
		Assume also the absorption condition
		\begin{equation}\label{eq:e-absorb}
			\bar e_k\le \frac{\eta}{8}\lambda_k \qquad\text{for some }\eta\in(0,1).
		\end{equation}
		Then
		\begin{equation}\label{eq:lambda-comp}
			c_\lambda\sqrt{M\norm{g_k}}\le \lambda_k \le C_\lambda\sqrt{M\norm{g_k}},
		\end{equation}
		where
		\[
		c_\lambda:=\frac{\sqrt{1-\alpha_d}}{\sqrt{1+\alpha_b}},
		\qquad
		C_\lambda:=\frac{1}{1-\eta/8}\cdot \frac{\sqrt{1+\alpha_d}}{\sqrt{1-\alpha_b}}.
		\]
	\end{lemma}
	
	\subsection{Descent and complexity from inexact regularized Newton}
	Define
	\[
	\Phi_k:=f(\bar x_k)-f_\star,\qquad g_k:=\nabla f(\bar x_k),\qquad \lambda_k:=\tilde{\bar\lambda}_k.
	\]
	
	\begin{lemma}[One-step descent]\label{lem:descent-inexact}
		Assume
		\[
		\norm{(\nabla^2 f(\bar x_k)+\lambda_k I)\bar s_k+g_k}\le \eta\,\lambda_k\norm{\bar s_k},
		\qquad
		L_2\norm{\bar s_k}\le \lambda_k,
		\]
		for some $\eta\in(0,1)$.
		Then
		\[
		f(\bar x_{k+1})\le f(\bar x_k) - \Big(1-\eta-\frac16\Big)\lambda_k\norm{\bar s_k}^2.
		\]
	\end{lemma}
	
	\begin{lemma}[Gradient control]\label{lem:grad-inexact}
		Under the conditions of Lemma~\ref{lem:descent-inexact},
		\[
		\norm{g_{k+1}}\le \Big(1+\eta+\frac12\Big)\lambda_k\norm{\bar s_k}.
		\]
	\end{lemma}
	
	Define steady/sharp sets:
	\[
	\mathcal I:=\left\{k\ge 0:\ \norm{g_{k+1}}\ge \frac14\norm{g_k}\right\},
	\qquad
	\mathcal S:=\left\{k\ge 0:\ \norm{g_{k+1}}< \frac14\norm{g_k}\right\}.
	\]
	
	\begin{lemma}[Steady steps yield a $3/2$ recursion]\label{lem:steady-rec}
		Assume Proposition~\ref{prop:inexact-rn} holds for all $k$ with some $\eta\le 1/12$ and that
		\eqref{eq:rel-bridge}--\eqref{eq:rel-disp} and \eqref{eq:e-absorb} hold for all $k$ (fixed $\alpha_b,\alpha_d\in(0,1/2]$).
		Then there exists $\tau>0$ such that for all $k\in\mathcal I$,
		\[
		\Phi_k-\Phi_{k+1}\ge \tau\,\Phi_k^{3/2}.
		\]
	\end{lemma}
	
	\begin{lemma}[Sharp steps are few]\label{lem:sharp-count}
		Let $\varepsilon>0$ and $K_\varepsilon:=\min\{k:\ \norm{g_k}\le \varepsilon\}$.
		Then
		\[
		|\mathcal S\cap\{0,1,\dots,K_\varepsilon-1\}|
		\le \left\lceil \log_4\frac{\norm{g_0}}{\varepsilon}\right\rceil.
		\]
	\end{lemma}
	
	\begin{lemma}[Solving the $3/2$ recursion]\label{lem:telescoping}
		If $\Phi_{k+1}\le \Phi_k-\tau \Phi_k^{3/2}$ for some $\tau>0$ and all $k$, then
		$\Phi_k\le \frac{4}{\tau^2(k+2)^2}$.
	\end{lemma}
	
	\begin{theorem}[Main complexity result]\label{thm:main}
		Assume Assumption~\ref{ass:standing}.
		Assume Proposition~\ref{prop:inexact-rn} holds for all $k$ with some $\eta\le 1/12$.
		Assume also \eqref{eq:rel-bridge}--\eqref{eq:rel-disp} and \eqref{eq:e-absorb} hold for all $k$ (fixed $\alpha_b,\alpha_d\in(0,1/2]$).
		Then:
		\begin{enumerate}[leftmargin=1.2em,itemsep=2pt]
			\item $\Phi_k=\cO(1/k^2)$.
			\item $\norm{g_k}=\cO(1/k)$.
			\item To reach $\norm{g_k}\le \varepsilon$, it suffices that $k=\cO(\varepsilon^{-1})$.
		\end{enumerate}
	\end{theorem}
	
	\subsection{Communication complexity under an explicit logarithmic schedule}\label{ssec:comm_schedule}
	We now remove the remaining circularity in the mixing discussion by fixing an explicit schedule for
	$(\tau_k,t_k)$ and proving that it enforces the accuracy conditions required by
	Proposition~\ref{prop:inexact-rn} and Theorem~\ref{thm:main} after a finite burn-in.
	
	\noindent\textbf{Logarithmic schedule.}
	Fix an exponent $p>2$ and define
	\begin{equation}\label{eq:log_schedule}
		\tau_k:=t_k:=\left\lceil \frac{p\,\log(k+2)}{-\log\rho}\right\rceil,\qquad k\ge 0.
	\end{equation}
	Then $\rho^{\tau_k}=\rho^{t_k}\le (k+2)^{-p}$ for all $k$.
	
	\noindent\textbf{Polynomial decay of dispersion drivers.}
	Define the (post-mixing) dispersions
	\[
	D_X(k):=\sqrt{\frac1N\sum_{i=1}^N\norm{x_{i,k}-\bar x_k}^2},\qquad
	D_G(k):=\sqrt{\frac1N\sum_{i=1}^N\norm{g_{i,k}-\bar g_k}^2},\qquad
	D_H(k):=\sqrt{\frac1N\sum_{i=1}^N\normF{H_{i,k}-\bar H_k}^2},
	\]
	and their pre-mixed counterparts $D_X(\tilde X_k),D_G(\tilde G_k),D_H(\tilde H_k)$.
	Appendix~\ref{app:dispersion} derives explicit recursions for $(D_X,D_G,D_H)$ and proves:
	
	\begin{proposition}[Dispersion decay under \eqref{eq:log_schedule}]\label{prop:dispersion_decay}
		Under Assumption~\ref{ass:standing}, there exist finite constants $C_X,C_G,C_H>0$ such that for all $k\ge 0$,
		\begin{align}
			D_X(\tilde X_k) &\le \frac{C_X}{(k+2)^{p-1}},\qquad
			D_G(\tilde G_k) &\le \frac{C_G}{(k+2)^{p-1}},\qquad
			D_H(\tilde H_k) &\le \frac{C_H}{(k+2)^{p-1}}.
		\end{align}
		Consequently, $\Delta^g_k\le D_G(\tilde G_k)$ and $\Delta^H_k\le D_H(\tilde H_k)$.
	\end{proposition}
	
	\begin{proposition}[Burn-in enforcement for a prescribed stationarity level]\label{prop:burnin_enforce}
		Fix $p>2$ and use the schedule~\eqref{eq:log_schedule}. Fix any accuracy level $\varepsilon>0$ satisfying
		\begin{equation}\label{eq:eps_floor}
			\varepsilon_{\mathrm{floor}}
			:=\max\left\{
			\frac{16\,\sigma_g}{\sqrt N},\ 
			\frac{61440\,\bar e_{\max}^2}{M},\
			\left(\frac{960\sqrt{5/3}\,M_H\,\sigma_g}{\sqrt{MN}}\right)^{\!\!2/3},\
			\left(\frac{1920\sqrt{5/3}\,M_H^2\,\bar e_{\max}}{M^{3/2}}\right)^{\!\!2/3}
			\right\}.
		\end{equation}
		Then there exists an integer $K_0(\varepsilon)\ge 0$ (explicitly given in Appendix~\ref{app:dispersion}) such that
		for every $k\ge K_0(\varepsilon)$ with $\norm{\nabla f(\bar x_k)}\ge \varepsilon$, the following three statements hold:
		\begin{enumerate}[leftmargin=1.2em,itemsep=2pt]
			\item the relative bridge and dispersion conditions \eqref{eq:rel-bridge}--\eqref{eq:rel-disp} hold with fixed constants
			$\alpha_b=\alpha_d=1/4$;
			\item the Hessian absorption condition \eqref{eq:e-absorb} holds with the fixed choice $\eta:=1/12$; moreover, the Hessian bridge accuracy in Proposition~\ref{prop:inexact-rn}(3) holds with the same $\eta$;
			\item the inexact regularized Newton conditions of Proposition~\ref{prop:inexact-rn} hold (with the same $\eta$).
		\end{enumerate}
		Consequently, on the index set $\{k\ge K_0(\varepsilon):\ \norm{\nabla f(\bar x_k)}\ge \varepsilon\}$,
		all hypotheses of Theorem~\ref{thm:main} are satisfied.
	\end{proposition}
	
	\begin{theorem}[End-to-end complexity under an explicit schedule]\label{thm:end2end}
		Assume Assumption~\ref{ass:standing}. Fix $p>2$ and run \textsc{DisGrem} with the schedule~\eqref{eq:log_schedule}.
		Let $\varepsilon>0$ satisfy \eqref{eq:eps_floor}.
		Define the first $\varepsilon$-stationary index
		\[
		K(\varepsilon):=\min\{k\ge 0:\ \norm{\nabla f(\bar x_k)}\le \varepsilon\}.
		\]
		Then there exist constants $C_\Phi,C_g,C_K>0$ (depending only on the problem and network parameters in Assumption~\ref{ass:standing} and on $p$) such that:
		\begin{enumerate}[leftmargin=1.2em,itemsep=2pt]
			\item (Functional decay above the floor) for all $0\le k\le K(\varepsilon)$,
			\[
			\Phi_k\le \frac{C_\Phi}{(k+2)^2};
			\]
			\item (Gradient decay above the floor) for all $0\le k\le K(\varepsilon)$,
			\[
			\norm{\nabla f(\bar x_k)}\le \frac{C_g}{k+2};
			\]
			\item (Iteration complexity) $K(\varepsilon)\le C_K\,\varepsilon^{-1}$.
		\end{enumerate}
	\end{theorem}
	\begin{theorem}[Total mixing rounds] \label{thm:comm}
		Assume the hypotheses of Theorem~\ref{thm:end2end} (in particular, the explicit schedule~\eqref{eq:log_schedule}).
		Let $K(\varepsilon):=\min\{k:\ \norm{\nabla f(\bar x_k)}\le \varepsilon\}$. Then
		\[
		\sum_{k=0}^{K(\varepsilon)-1}(\tau_k+t_k)
		=\cO\bigl(\varepsilon^{-1}\log(\varepsilon^{-1})\bigr)
		\]
		for any $\varepsilon$ above the accuracy floor.
	\end{theorem}
	
	
	\section{Numerical Experiments}\label{sec:experiments}
	%%%%%%%%%%%%%%%%%%%%%%%%%%%%%%%%%%%%%%%%%%%%%%%%%%%%%%%%%%%%
	
	We evaluate \textsc{DisGrem} and four variants against a comprehensive suite of first- and
	second-order decentralized baselines across 15 benchmark objectives spanning
	well-conditioned and ill-conditioned convex problems, real-data tasks, and non-convex landscapes.
	
	% ─────────────────────────────────────────────────────────────────────
	\subsection{Experimental setup}\label{ssec:setup}
	% ─────────────────────────────────────────────────────────────────────
	
	\noindent\textbf{Network.}
	We use $N=10$ agents on a random geometric graph with connection radius $r_0=0.5$
	in the unit square (regenerated per trial).
	The Metropolis--Hastings doubly stochastic matrix $W$ is formed from the resulting
	adjacency; the spectral norm $\rho=\normtwo{W-\J}\approx 0.92$ (i.e.\ the
	second-largest absolute eigenvalue of $W$), representing a moderately sparse,
	weakly connected topology with spectral gap $1-\rho\approx 0.08$.
	All algorithms perform $N_c=3$ consensus rounds per outer iteration ($\tau_k=t_k=3$).
	
	\noindent\textbf{Algorithms.}
	We compare 10 algorithms in three groups.
	
	\noindent\textit{Proposed (\textsc{DisGrem} family)}:
	\textbf{DisGrem} (Algorithm~\ref{alg:disgrem}), its compressed-Hessian variant \textbf{CeDisGrem}
	(top-$r$ spectral compression of Hessian corrections, Section~\ref{ssec:variants}),
	the adaptive-$M$ variant \textbf{AdaDisGrem}, and its compressed version \textbf{CeAdaDisGrem}.
	The regularization parameter is set to $M = M_{\mathrm{fac}} \cdot \hat L$
	where $\hat L$ estimates the gradient Lipschitz constant
	(with $M_{\mathrm{fac}}$ studied in Section~\ref{ssec:param}).
	
	\noindent\textit{First-order baselines}:
	\textbf{EXTRA}~\citep{shi2015extra} (exact first-order decentralized method),
	\textbf{DIGing}~\citep{nedic2017diging} (gradient-tracking DGD),
	and \textbf{DGD+GT} (standard DGD enhanced with gradient tracking~\citep{nedic2009subgradient}).
	EXTRA uses its canonical step size $\alpha\approx 1/(2L)$;
	DIGing uses $\alpha\approx 0.5/L$; DGD+GT uses $\alpha=0.2/L$.
	
	\noindent\textit{Second-order baselines}:
	\textbf{DQM}~\citep{eisen2017dqn} (decentralized quasi-Newton),
	\textbf{ESOM} (exact second-order method with network Newton--style Hessian~\citep{mokhtari2015networknewton_arxiv}),
	and \textbf{NT} (Newton Tracking~\citep{mokhtari2020newtontracking}).
	All second-order baselines use their published default stepsizes.
	
	\noindent\textbf{Test objectives.}
	We use 15 objectives in two groups (dimension $d=10$ unless noted).
	\textit{Convex} ($9$ problems): (i)~Ridge regression (\textsc{Ridge}); (ii)~ill-conditioned
	quadratic (\textsc{QuadBad}, $\kappa=10^3$); (iii)~log-sum-exp (\textsc{LSE}); (iv)~Huber loss;
	(v)~Tsallis entmax (\textsc{EntMax}); (vi)~log-linear (\textsc{LinLog});
	(vii)~pseudo-Huber + quadratic (\textsc{AbsQuad}, smoothing scale $\delta=0.5$);
	(viii)~$\ell_2$-regularized logistic regression on a real LibSVM dataset
	(\textsc{LR-real}, $d=22$); (ix)~least-squares regression on the scikit-learn
	Diabetes dataset (\textsc{LS-real}).
	\textit{Non-convex} ($6$ problems): Rosenbrock, Styblinski--Tang, Ackley,
	SinQuad, CosMix, and a non-convex variant of logistic regression (\textsc{LR-ncvr}).
	Reference optima are computed by high-precision LBFGS-B (tolerance $10^{-15}$) with multiple restarts.
	
	\noindent\textbf{Stopping criterion and metrics.}
	All algorithms run to a maximum of $K_{\max}$ iterations (500--2000, depending on problem
	difficulty) or until the \emph{combo} criterion $\|\bar g_k\|+\mathrm{cons}_k < 10^{-12}$
	is satisfied, where $\bar g_k$ is the tracked gradient average and $\mathrm{cons}_k$ is the
	per-agent disagreement.
	We report: (a)~\emph{Steps}: outer iterations to stopping;
	(b)~\emph{combo}: $\min_k(\|\bar g_k\|+\mathrm{cons}_k)$;
	(c)~\emph{relF}: $\min_k |f(\bar x_k)-f^\star|/|f(\bar x_0)-f^\star|$; and
	(d)~\emph{Comm}: average total communication in MB.
	Analytical gradients and Hessians are used whenever available.

	\noindent\textbf{Fairness and reproducibility protocol.}
	To ensure a fair comparison, all methods share the same network realization per trial,
	initialization policy, stopping tolerance, and communication accounting rule
	(per-agent MB over all neighbor exchanges).
	Every method is run from identical random seeds for each trial.
	Unless explicitly varied in a sensitivity study, hyperparameters follow either
	(i) published default recommendations or (ii) a single global rule (e.g., $M=M_{\mathrm{fac}}\hat L$)
	applied uniformly across problems.
	We emphasize that no per-problem hand-retuning is performed for baselines beyond
	their standard stable range.
	
	% ─────────────────────────────────────────────────────────────────────
	\subsection{Convergence on convex benchmarks}\label{ssec:convex}
	% ─────────────────────────────────────────────────────────────────────
	
	Table~\ref{tab:main} summarises results on seven representative convex objectives.
	Figure~\ref{fig:multiobj} shows iteration-vs-combo convergence curves across
	all 15 objectives (9 convex + 6 non-convex).
	
	\begin{table}[t]
		\centering
		\caption{Convergence on convex benchmarks ($N=10$, $d=10$ or $22$).
			\textbf{Steps}: outer iterations at stopping (or $K_{\max}$ if not converged;
			$K_{\max}=500$ for Ridge/Huber/LinLog/LS-real, $800$ for LSE, $2000$ for EntMax/LR-real).
			\textbf{combo}: $\min_k(\|\bar g_k\|+\mathrm{cons}_k)$; $\checkmark$ denotes combo $<10^{-12}$.
			\textbf{MB}: per-agent total communication in megabytes.
			Bold rows = proposed methods.}
		\label{tab:main}
		\setlength{\tabcolsep}{4pt}
		\resizebox{\textwidth}{!}{%
			\begin{tabular}{l ccc ccc ccc ccc ccc ccc ccc}
				\toprule
				& \multicolumn{3}{c}{\textsc{Ridge}}
				& \multicolumn{3}{c}{\textsc{LSE}}
				& \multicolumn{3}{c}{\textsc{Huber}}
				& \multicolumn{3}{c}{\textsc{EntMax}}
				& \multicolumn{3}{c}{\textsc{LinLog}}
				& \multicolumn{3}{c}{\textsc{LR-real}}
				& \multicolumn{3}{c}{\textsc{LS-real}} \\
				\cmidrule(lr){2-4}\cmidrule(lr){5-7}\cmidrule(lr){8-10}%
				\cmidrule(lr){11-13}\cmidrule(lr){14-16}\cmidrule(lr){17-19}\cmidrule(lr){20-22}
				Algorithm & Steps & combo & MB & Steps & combo & MB & Steps & combo & MB
				& Steps & combo & MB & Steps & combo & MB & Steps & combo & MB
				& Steps & combo & MB \\
				\midrule
				\multicolumn{22}{l}{\textit{Proposed}} \\
				\textbf{DisGrem}      & 416 & $\checkmark$ & 4.3
				& 441 & $\checkmark$ & 4.6
				& 378 & $\checkmark$ & 3.9
				& 689 & $\checkmark$ & 7.1
				& 271 & $\checkmark$ & 2.8
				& 1273 & $\checkmark$ & 59.7
				& 500 & $8\times10^{-5}$ & 5.2 \\
				\textbf{CeDisGrem}    & 412 & $\checkmark$ & 2.9
				& 429 & $\checkmark$ & 3.0
				& 428 & $\checkmark$ & 3.0
				& 766 & $\checkmark$ & 5.3
				& 329 & $\checkmark$ & 2.3
				& 1578 & $\checkmark$ & 49.4
				& 500 & $10^{-8}$    & 3.5 \\
				\textbf{AdaDisGrem}   & 423 & $\checkmark$ & 4.4
				& 402 & $\checkmark$ & 4.2
				& 386 & $\checkmark$ & 4.0
				& 704 & $\checkmark$ & 7.3
				& 275 & $\checkmark$ & 2.8
				& 1343 & $\checkmark$ & 62.9
				& 500 & $9\times10^{-5}$ & 5.2 \\
				\textbf{CeAdaDisGrem} & 421 & $\checkmark$ & 2.9
				& 376 & $\checkmark$ & 2.6
				& 439 & $\checkmark$ & 3.0
				& 779 & $\checkmark$ & 5.4
				& 337 & $\checkmark$ & 2.3
				& 1645 & $\checkmark$ & 51.5
				& 500 & $8\times10^{-8}$ & 3.5 \\
				\midrule
				\multicolumn{22}{l}{\textit{First-order baselines}} \\
				EXTRA  & 500  & $4\times10^{-6}$ & 1.6
				& 800  & $7\times10^{-8}$ & 2.5
				& 500  & $1\times10^{-2}$ & 1.6
				& 2000 & $3\times10^{-4}$ & 6.3
				& 500  & $2\times10^{-2}$ & 1.6
				& 2000 & $2.1$            & 15.1
				& 500  & $1\times10^{-3}$ & 1.6 \\
				DIGing & 500  & $5\times10^{-1}$ & 3.4
				& 800  & $6\times10^{-5}$ & 5.5
				& 500  & $1\times10^{-1}$ & 3.4
				& 2000 & $4\times10^{-3}$ & 13.7
				& 500  & $5\times10^{-1}$ & 3.4
				& 2000 & $3.9\times10^{-1}$ & 62.4
				& 500  & $6\times10^{-3}$ & 3.4 \\
				DGD+GT & 500  & $6\times10^{-1}$ & 3.4
				& 800  & $4\times10^{-3}$ & 5.5
				& 500  & $7\times10^{-1}$ & 3.4
				& 2000 & $4\times10^{-2}$ & 13.7
				& 500  & $3.4$            & 3.4
				& 2000 & $4\times10^{-11}$ & 62.5
				& 500  & $3\times10^{-2}$ & 3.4 \\
				\midrule
				\multicolumn{22}{l}{\textit{Second-order baselines}} \\
				DQM   & 500  & $4\times10^{-2}$ & 4.0
				& 800  & $6\times10^{-2}$ & 6.4
				& 500  & $5\times10^{-1}$ & 4.0
				& 2000 & $3\times10^{-1}$ & 11.2
				& 500  & $9\times10^{-1}$ & 4.0
				& 2000 & $1\times10^{-2}$ & 72.9
				& 500  & $3\times10^{-2}$ & 4.0 \\
				ESOM  & 500  & $4\times10^{-2}$ & 10.3
				& 800  & $7\times10^{-2}$ & 16.5
				& 500  & $6.2$            & 10.3
				& 2000 & $2.0$            & 37.5
				& 500  & $897$            & 10.3
				& 2000 & $9.9\times10^{-3}$ & 187.4
				& 500  & $3\times10^{-2}$ & 10.3 \\
				NT    & 500  & $6\times10^{-2}$ & 1.2
				& 800  & $1.5$            & 1.8
				& 500  & $3.8$            & 1.2
				& 2000 & $2.2$            & 4.6
				& 500  & $3.3$            & 1.2
				& 2000 & $7.6\times10^{-1}$ & 20.8
				& 500  & $6.6\times10^{-1}$ & 1.2 \\
				\bottomrule
		\end{tabular}}
	\end{table}
	
	\begin{figure}[t]
		\centering
		\begin{subfigure}[b]{0.49\textwidth}
			\safeincludegraphics[width=\textwidth]{multiobj_steps_vs_combo}
			\caption{Steps vs.\ combo ($\|\bar g\|+\mathrm{cons}$)}
		\end{subfigure}
		\hfill
		\begin{subfigure}[b]{0.49\textwidth}
			\safeincludegraphics[width=\textwidth]{multiobj_commCost_vs_relF}
			\caption{Communication (MB) vs.\ relF}
		\end{subfigure}
		\caption{Convergence curves across all 15 benchmarks (9 convex + 6 non-convex).
			Each panel shows one function; shared legend at bottom.
			Proposed algorithms (\textbf{bold lines}) consistently reach combo~$<10^{-12}$
			on the convex suite, while first-order baselines often stagnate on weakly connected graphs.}
		\label{fig:multiobj}
	\end{figure}
	
	\noindent\textbf{Key observations.}
	\begin{enumerate}[leftmargin=1.4em,itemsep=2pt]
		\item \textbf{Proposed algorithms converge to machine precision on all well-conditioned convex problems.}
		All four variants achieve combo~$<10^{-12}$ on Ridge (412--423 steps),
		LSE (376--441), Huber (378--439), EntMax (689--779), and LinLog (271--337).
		On the real-data task LR-real ($d=22$), all four converge within 1273--1645 steps.
		On LS-real, compressed variants reach combo~$\approx10^{-8}$ (within $K_{\max}=500$),
		while uncompressed variants are slightly slower to converge on this mildly ill-conditioned
		problem, illustrating the benefit of Hessian compression for larger residuals.
		
		\item \textbf{EXTRA is the strongest first-order baseline yet substantially slower.}
		On Ridge, EXTRA reaches combo~$\approx 4\times10^{-6}$ after $K_{\max}=500$ steps---six
		orders of magnitude above machine precision---while our algorithms converge fully in 412--423 steps.
		On EntMax, EXTRA needs all 2000 steps and still achieves only combo~$\approx 3\times10^{-4}$;
		our methods converge in 689--779 steps.
		Although EXTRA's per-iteration communication (two $d$-vectors per consensus round) is
		lower than ours (one vector plus one Hessian matrix), the large iteration disparity
		more than offsets this advantage: EXTRA still fails to satisfy the target stopping criterion
		on harder tasks even after exhausting the full communication budget.
		
		\item \textbf{DIGing and DGD+GT struggle on weakly connected networks.}
		Both first-order gradient-tracking methods fail to make significant progress on Ridge and LinLog
		within 500 iterations (combo $\approx 0.5$--$3.4$).
		The underlying cause is the small spectral gap ($1-\rho\approx0.08$):
		first-order consensus-gradient methods require $\cO((1-\rho)^{-2})$ additional iterations
		to suppress disagreement~\citep{nedic2009subgradient},
		whereas the Hessian-regularized Newton step in \textsc{DisGrem} produces far larger
		per-step progress, rendering the algorithm less sensitive to network connectivity.
		
		\item \textbf{Second-order baselines are inconsistent across function classes.}
		DQM and ESOM perform reasonably on Ridge but diverge or stagnate on LinLog, Huber, and EntMax.
		Newton Tracking (NT) struggles on LSE and LR-real due to Hessian rank deficiency
		and ill-conditioning.
		Our methods are the only ones achieving combo~$<10^{-12}$ consistently across all
		well-conditioned convex problems.
		
		\item \textbf{Communication efficiency.}
		The compressed variants CeDisGrem and CeAdaDisGrem reduce communication by
		$30$--$50\%$ relative to their uncompressed counterparts (e.g.\ Ridge: 2.9\,MB vs.\ 4.3\,MB)
		with negligible impact on iteration count, confirming that Hessian compression is effective.
	\end{enumerate}

	\begin{table}[t]
		\centering
		\caption{Mechanism-level ablation on representative convex tasks.
			$\Delta$Steps and $\Delta$MB are relative changes versus \textbf{DisGrem}
			(negative is better).}
		\label{tab:ablation}
		\begin{tabular}{lcccc}
			\toprule
			Task & CeDisGrem $\Delta$Steps & CeDisGrem $\Delta$MB & AdaDisGrem $\Delta$Steps & AdaDisGrem $\Delta$MB \\
			\midrule
			Ridge   & $-1.0\%$  & $-32.6\%$ & $+1.7\%$ & $+2.3\%$ \\
			LSE     & $-2.7\%$  & $-34.8\%$ & $-8.8\%$ & $-8.7\%$ \\
			EntMax  & $+11.2\%$ & $-25.4\%$ & $+2.2\%$ & $+2.8\%$ \\
			LR-real & $+23.9\%$ & $-17.3\%$ & $+5.5\%$ & $+5.4\%$ \\
			\bottomrule
		\end{tabular}
	\end{table}
	
	\noindent\textbf{Non-convex problems.}
	Although our convergence theory (Theorem~\ref{thm:end2end}) requires convexity,
	we include non-convex benchmarks to test practical behavior.
	Figure~\ref{fig:multiobj} includes two non-convex functions (Rosenbrock, Styblinski--Tang)
	where DisGrem family curves are plotted.
	On Rosenbrock, DisGrem and AdaDisGrem converge in 107 and 156 steps to
	combo~$<10^{-13}$ and relF~$\approx 10^{-20}$.
	On Styblinski--Tang (a standard multimodal benchmark), all four variants reach
	combo~$<10^{-9}$ within 1000 steps; first-order methods fail entirely
	(EXTRA: combo~$=18.5$, DGD+GT: combo~$>100$).
	On SinQuad and CosMix the DisGrem family reliably finds a stationary point
	(combo~$<10^{-11}$), while the first-order baselines stagnate.
	These results suggest that the regularized Newton structure provides sufficient
	implicit regularization to navigate non-convex landscapes in practice,
	though formal non-convex guarantees are left for future work.
	
	% ─────────────────────────────────────────────────────────────────────
	\subsection{Monte-Carlo robustness study}\label{ssec:robust}
	% ─────────────────────────────────────────────────────────────────────
	
	To assess sensitivity to initialisation, we run $n_s=10$ independent trials on the
	Ridge problem with different random seeds (both initial point and network topology).
	A trial is declared \emph{successful} if relF $<10^{-5}$ is achieved within 1000 iterations.
	Table~\ref{tab:robust} reports success rate, average steps (over successful runs),
	and average wall-clock time; Figure~\ref{fig:robust} visualises the results.
	
	\begin{table}[h]
		\centering
		\caption{Monte-Carlo robustness on Ridge ($N=10$, $d=10$, $n_s=10$ trials).
			\textbf{Succ}: success rate (\%). \textbf{Avg-K}: mean steps over successful runs.
			\textbf{Avg-t}: mean wall-clock time over successful runs (seconds).}
		\label{tab:robust}
		\begin{tabular}{lccc}
			\toprule
			Algorithm & Succ (\%) & Avg-K & Avg-t (s) \\
			\midrule
			\textbf{DisGrem}    & \textbf{100} & \textbf{30.6} & 0.07 \\
			\textbf{CeDisGrem}  & \textbf{100} & \textbf{22.3} & 0.05 \\
			EXTRA               & 100          & 162.6         & 1.46 \\
			DIGing              & 20           & 339.5         & 3.22 \\
			DGD+GT              & 0            & ---           & --- \\
			\bottomrule
		\end{tabular}
	\end{table}
	
	\begin{figure}[h]
		\centering
		\safeincludegraphics[width=0.55\textwidth]{robust_ridge}
		\caption{Robustness study on Ridge ($n_s=10$ random starts).
			Left: success rate per algorithm; right: box plot of steps for successful runs.}
		\label{fig:robust}
	\end{figure}
	
	The results are striking.
	\textbf{DisGrem and CeDisGrem achieve 100\% success with a mean of only 30.6 and 22.3 steps},
	roughly $5\times$ and $7\times$ fewer than EXTRA (162.6 steps, also 100\% success).
	DIGing succeeds in only 2 out of 10 runs, and DGD+GT \emph{never} converges within budget.
	This illustrates that the weak spectral gap ($1-\rho\approx0.08$) is catastrophic for
	first-order gradient-tracking methods but is naturally handled by our second-order scheme.
	To quantify uncertainty under finite trials, the 95\% Clopper--Pearson interval for success rate is
	$[69.1\%,100\%]$ for DisGrem/CeDisGrem/EXTRA (10/10), $[2.5\%,55.6\%]$ for DIGing (2/10),
	and $[0\%,30.8\%]$ for DGD+GT (0/10), confirming a statistically meaningful separation.
	
	% ─────────────────────────────────────────────────────────────────────
	\subsection{Scalability}\label{ssec:scale}
	% ─────────────────────────────────────────────────────────────────────
	
	We examine how each algorithm scales as (a) the number of agents $N$ grows and
	(b) the problem dimension $d$ grows, using Ridge as the benchmark.
	
	\noindent\textbf{Varying $N$ ($d=10$ fixed).}
	Figure~\ref{fig:scale_N} and Table~\ref{tab:scale_N} report steps to convergence for
	$N\in\{5,10,20,50\}$.
	The DisGrem family converges in 26--353 steps across all sizes, with step counts
	\emph{decreasing} as $N$ grows (90--120 steps at $N=50$)---a consequence of more agents
	providing richer curvature information.
	EXTRA converges in 833--868 steps regardless of $N$, showing much weaker dependence on agent count.
	DIGing and DGD+GT do not converge within 1500 steps for any $N$.
	Across all tested $N$, the proposed methods retain a large margin over first-order methods,
	indicating that increasing agent count does not destabilize second-order tracking dynamics.
	
	\begin{table}[h]
		\centering
		\caption{Steps to convergence (combo $<10^{-10}$) vs.\ number of agents $N$
			($d=10$, Ridge). ``NC'' = not converged within 1500 steps.}
		\label{tab:scale_N}
		\begin{tabular}{lrrrr}
			\toprule
			Algorithm      & $N=5$ & $N=10$ & $N=20$ & $N=50$ \\
			\midrule
			\textbf{DisGrem}    & 67  & 307 & 136 & 90  \\
			\textbf{CeDisGrem}  & 26  & 353 & 149 & 120 \\
			\textbf{AdaDisGrem} & 70  & 303 & 139 & 94  \\
			\textbf{CeAdaDisGrem} & 29 & 351 & 149 & 121 \\
			EXTRA               & 846 & 868 & 852 & 833 \\
			DIGing              & NC  & NC  & NC  & NC  \\
			DGD+GT              & NC  & NC  & NC  & NC  \\
			\bottomrule
		\end{tabular}
	\end{table}
	
	\noindent\textbf{Varying $d$ ($N=10$ fixed).}
	Figure~\ref{fig:scale_d} shows steps and wall-clock time for $d\in\{10,30,50,100\}$.
	DisGrem and CeDisGrem converge for all values of $d$:
	step counts remain in the range 207--738, and CeDisGrem consistently converges faster
	than DisGrem for larger $d$ (685 vs.\ 738 at $d=100$) with substantially lower
	communication (108\,MB vs.\ 172\,MB), confirming that Hessian compression becomes
	increasingly beneficial as dimension grows.
	EXTRA, DIGing, and DGD+GT fail to converge for $d\ge30$.
	This trend is consistent with the communication model in Section~\ref{ssec:comm}:
	the Hessian payload scales quadratically with $d$, so low-rank transmission brings larger gains
	in high dimensions.
	
	\begin{figure}[t]
		\centering
		\begin{subfigure}[b]{0.49\textwidth}
			\safeincludegraphics[width=\textwidth]{scalability_N_steps_ridge}
			\caption{Steps vs.\ $N$ ($d=10$ fixed)}
			\label{fig:scale_N}
		\end{subfigure}
		\hfill
		\begin{subfigure}[b]{0.49\textwidth}
			\safeincludegraphics[width=\textwidth]{scalability_d_steps_ridge}
			\caption{Steps vs.\ $d$ ($N=10$ fixed)}
			\label{fig:scale_d}
		\end{subfigure}
		\caption{Scalability of DisGrem family (solid) vs.\ EXTRA (dashed) on Ridge.
			``NC'' entries (DIGing, DGD+GT) are omitted from the plot as they never converge.
			Note log scale on $y$-axis.}
		\label{fig:scalability}
	\end{figure}
	
	% ─────────────────────────────────────────────────────────────────────
	\subsection{Sensitivity to the regularisation parameter $M$}\label{ssec:param}
	% ─────────────────────────────────────────────────────────────────────
	
	We sweep $M_{\mathrm{fac}}\in\{0.05,\,0.1,\,0.5,\,1.0,\,3.0\}\times L_{\max}$
	for four representative algorithms on Ridge ($N=10$, $d=10$).
	Figure~\ref{fig:param} shows steps-vs-combo curves for DisGrem and EXTRA;
	Figure~\ref{fig:param_ce} shows CeDisGrem and DIGing.
	
	\begin{figure}[t]
		\centering
		\begin{subfigure}[b]{0.49\textwidth}
			\safeincludegraphics[width=\textwidth]{DisGrem_steps_vs_combo}
			\caption{\textbf{DisGrem}: all $M_{\mathrm{fac}}$ converge}
		\end{subfigure}
		\hfill
		\begin{subfigure}[b]{0.49\textwidth}
			\safeincludegraphics[width=\textwidth]{EXTRA_steps_vs_combo}
			\caption{EXTRA: only the theoretically tuned step size converges; aggressive steps diverge}
		\end{subfigure}
		\caption{Parameter sensitivity on Ridge.
			For DisGrem, each curve corresponds to one value of $M_{\mathrm{fac}}$;
			for EXTRA, each curve corresponds to one step-size multiplier $c_\alpha\in\{0.2,0.5,1.0,2.0,5.0\}$
			times the reference step $1/(2L)$.
			\textbf{DisGrem converges for all five $M_{\mathrm{fac}}$ settings}; the rate
			improves as $M_{\mathrm{fac}}$ grows from $0.05$ to $1.0$, consistent with the
			theoretical prediction that larger $M$ expands the admissible heterogeneity regime.
			EXTRA's safe range is narrow: only the reference step size converges; doubling or halving leads to stagnation or divergence.}
		\label{fig:param}
	\end{figure}
	
	\begin{figure}[t]
		\centering
		\begin{subfigure}[b]{0.49\textwidth}
			\safeincludegraphics[width=\textwidth]{CeDisGrem_steps_vs_combo}
			\caption{\textbf{CeDisGrem}: all $M_{\mathrm{fac}}$ converge}
		\end{subfigure}
		\hfill
		\begin{subfigure}[b]{0.49\textwidth}
			\safeincludegraphics[width=\textwidth]{DIGing_steps_vs_combo}
			\caption{DIGing: sensitive to step size}
		\end{subfigure}
		\caption{Parameter sensitivity (continued).
			For CeDisGrem, each curve is one $M_{\mathrm{fac}}$ value; for DIGing, each curve is one
			step-size multiplier.
			\textbf{CeDisGrem} mirrors DisGrem's robustness while reducing communication by 30--50\%.
			\textbf{DIGing}'s convergence is highly sensitive to step size: three of five configurations
			fail to converge within budget.}
		\label{fig:param_ce}
	\end{figure}
	
	\noindent\textbf{Discussion.}
	All four DisGrem-family algorithms converge for every tested $M_{\mathrm{fac}}$,
	with the fastest convergence around $M_{\mathrm{fac}}\in[0.5,1.0]$.
	This robustness is consistent with the theoretical analysis: $M$ acts as a
	pure amplification constant that enlarges the admissible heterogeneity regime
	(the discussion in Section~\ref{sec:theory}) and does not appear as an explicit stepsize.
	In contrast, EXTRA and DIGing are highly sensitive to their respective step-size
	parameters: both show divergence or prolonged stagnation when the step size deviates
	from its theoretically prescribed value by a factor of 2--5,
	confirming that their convergence guarantees are inherently tied to exact step-size tuning.

	% ─────────────────────────────────────────────────────────────────────
	\subsection{What the experiments validate (and what they do not)}\label{ssec:exp_validity}
	% ─────────────────────────────────────────────────────────────────────

	\noindent\textbf{Validated claims.}
	The empirical evidence supports three central claims of this paper:
	(i) robust high-accuracy convergence of DisGrem-family methods on weakly connected networks,
	(ii) communication--accuracy benefits of Hessian compression, and
	(iii) broad parameter stability with respect to $M$.

	\noindent\textbf{Current limitations.}
	Two limitations remain: first, the strongest robustness analysis is currently reported on Ridge only;
	second, all experiments are deterministic/full-batch and do not cover stochastic gradients/Hessians.
	These are important extensions for a camera-ready large-scale ML benchmark suite.

	\noindent\textbf{Immediate next experimental add-ons.}
	To further strengthen the paper for top-journal review, the highest-priority add-ons are:
	(1) robustness trials on at least two additional objectives (one convex ill-conditioned, one non-convex),
	(2) a larger network sweep including denser and sparser graphs at fixed $N$,
	and (3) one real high-dimensional dataset ($d\gg 100$) to stress communication scaling.
	
	%%%%%%%%%%%%%%%%%%%%%%%%%%%%%%%%%%%%%%%%%%%%%%%%%%%%%%%%%%%%
	\section{Conclusion}
	%%%%%%%%%%%%%%%%%%%%%%%%%%%%%%%%%%%%%%%%%%%%%%%%%%%%%%%%%%%%
	
	\noindent\textbf{Summary.}
	We proposed \textsc{DisGrem}, a fully decentralized, stepsize-free second-order optimization
	method that combines gradient tracking, Hessian tracking with PSD projection,
	and two-stage mixing.
	The key analytical contribution is a \emph{reference-step framework} that resolves the
	non-commutativity between network averaging and Newton-system inversion,
	reducing the average iterate's dynamics to an \emph{inexact regularized Newton} recursion.
	By quantifying Hessian heterogeneity via the bounded constant $\sigma_H$ on the relevant level set
	and choosing $\lambda_{i,k}=\sqrt{M\|\tilde g_{i,k}\|}$ with $M\ge L_2$,
	we prove:
	\begin{itemize}[leftmargin=1.2em,itemsep=2pt]
		\item $\Phi_k = \cO(1/k^2)$ functional decay and $\cO(\varepsilon^{-1})$ iteration complexity,
		matching the centralized regularized Newton rate;
		\item $\cO(\varepsilon^{-1}\log(\varepsilon^{-1}))$ total neighbor-communication rounds
		under a logarithmic accuracy schedule.
	\end{itemize}
	
	\noindent\textbf{Experimental findings.}
	Extensive experiments across 15 benchmark objectives substantiate the theory:
	\begin{itemize}[leftmargin=1.2em,itemsep=2pt]
		\item \textsc{DisGrem} and its variants are the \emph{only} algorithms achieving
		combo~$<10^{-12}$ (gradient norm plus consensus residual) on all tested convex problems,
		with $5$--$8\times$ fewer iterations than EXTRA and orders-of-magnitude better accuracy
		than DQM, ESOM, and Newton Tracking.
		\item On weakly connected networks ($\rho\approx 0.92$), first-order gradient-tracking methods
		(EXTRA, DIGing, DGD+GT) stagnate or oscillate, while \textsc{DisGrem} converges
		reliably due to its second-order steps effectively damping network-induced disagreement.
		\item In 10-trial Monte-Carlo robustness tests, \textsc{DisGrem} achieves 100\% success
		with a mean of only 30.6 iterations (CeDisGrem: 22.3 iterations);
		EXTRA also succeeds but needs 162.6 iterations on average ($5$--$7\times$ more);
		and DIGing/DGD+GT fail in $\ge80\%$ of trials.
		\item \textsc{CeDisGrem}'s Hessian compression reduces communication by 30--50\% with
		negligible iteration overhead, becoming increasingly advantageous as dimension $d$ grows.
		\item Parameter sensitivity: all DisGrem variants converge for all tested scaling factors
		$M_\mathrm{fac}\in[0.05, 3.0]\times L_{\max}$, confirming the theory's prediction that
		$M$ acts as a pure amplifier rather than a sensitive stepsize.
	\end{itemize}
	
	\noindent\textbf{Limitations and future work.}
	Several directions merit further investigation:
	\begin{itemize}[leftmargin=1.2em,itemsep=2pt]
		\item \textbf{Heterogeneity dependence and accuracy floors.}
		The constants scale with the gradient heterogeneity $\sigma_g$ and curvature heterogeneity $\sigma_H$.
		When $\sigma_g>0$, an unavoidable stationarity floor appears (cf.\ \eqref{eq:eps_floor}).
		Designing heterogeneity-aware regularization or variance-reduction mechanisms to reduce these constants is an important direction.
		\item \textbf{Stochastic and mini-batch extensions.}
		The current analysis assumes exact local gradients and Hessians.
		Extending to stochastic approximations (e.g.\ via variance reduction) is important
		for large-scale machine learning applications.
		\item \textbf{Time-varying and directed graphs.}
		The mixing matrix $W$ is assumed symmetric and time-invariant.
		Handling directed graphs or dropping links would require a more delicate consensus
		analysis and may affect the $\cO(\varepsilon^{-1}\log(\varepsilon^{-1}))$ bound.
		\item \textbf{Lower bounds.}
		It is an open question whether $\cO(\varepsilon^{-1})$ iterations and
		$\cO(\varepsilon^{-1}\log(\varepsilon^{-1}))$ communication rounds are optimal for
		decentralized second-order methods under the stated assumptions.
	\end{itemize}
	
	%%%%%%%%%%%%%%%%%%%%%%%%%%%%%%%%%%%%%%%%%%%%%%%%%%%%%%%%%%%%
	\bibliographystyle{plainnat}
	\begin{thebibliography}{99}
		
		\bibitem[Cartis et~al.(2011)Cartis, Gould, and Toint]{cartis2011arc}
		C.~Cartis, N.~I.~M. Gould, and Ph.~L. Toint.
		\newblock Adaptive cubic regularisation methods for unconstrained optimization. Part~I: motivation, convergence and numerical results.
		\newblock \emph{Mathematical Programming}, 127:245--295, 2011.
		
		\bibitem[Eisen et~al.(2017)Eisen, Mokhtari, and Ribeiro]{eisen2017dqn}
		M.~Eisen, A.~Mokhtari, and A.~Ribeiro.
		\newblock Decentralized quasi-{N}ewton methods.
		\newblock \emph{IEEE Transactions on Signal Processing}, 65(10):2613--2628, 2017.
		
		\bibitem[Jakoveti{\'c} et~al.(2014)Jakoveti{\'c}, Xavier, and Moura]{jakovetic2014fast}
		D.~Jakoveti{\'c}, J.~Xavier, and J.~M.~F. Moura.
		\newblock Fast distributed gradient methods.
		\newblock \emph{IEEE Transactions on Automatic Control}, 59(5):1131--1146, 2014.
		
		\bibitem[Mishchenko(2023)]{mishchenko2023rn}
		K.~Mishchenko.
		\newblock Regularized {N}ewton method with global convergence and complexity bounds.
		\newblock \emph{arXiv preprint arXiv:2112.02089v3}, 2023.
		
		\bibitem[Mokhtari et~al.(2015)Mokhtari, Ling, and Ribeiro]{mokhtari2015networknewton_arxiv}
		A.~Mokhtari, Q.~Ling, and A.~Ribeiro.
		\newblock Network {N}ewton---part {I}: algorithm and convergence.
		\newblock \emph{arXiv preprint arXiv:1504.06017}, 2015.
		
		\bibitem[Mokhtari et~al.(2020)Mokhtari, Koppel, and Ribeiro]{mokhtari2020newtontracking}
		A.~Mokhtari, A.~Koppel, and A.~Ribeiro.
		\newblock Newton tracking: a decentralized method for exact convex optimization.
		\newblock \emph{IEEE Transactions on Signal Processing}, 68:572--587, 2020.
		
		\bibitem[Nedi{\'c} and Ozdaglar(2009)]{nedic2009subgradient}
		A.~Nedi{\'c} and A.~Ozdaglar.
		\newblock Distributed subgradient methods for multi-agent optimization.
		\newblock \emph{IEEE Transactions on Automatic Control}, 54(1):48--61, 2009.
		
		\bibitem[Nedi{\'c} et~al.(2017)Nedi{\'c}, Olshevsky, and Shi]{nedic2017diging}
		A.~Nedi{\'c}, A.~Olshevsky, and W.~Shi.
		\newblock Achieving geometric convergence for distributed optimization over time-varying graphs.
		\newblock \emph{SIAM Journal on Optimization}, 27(4):2597--2633, 2017.
		
		\bibitem[Nesterov and Polyak(2006)]{nesterov2006cubic}
		Y.~Nesterov and B.~T. Polyak.
		\newblock Cubic regularization of {N}ewton method and its global performance.
		\newblock \emph{Mathematical Programming}, 108:177--205, 2006.
		
		\bibitem[Shi et~al.(2015)Shi, Ling, Wu, and Yin]{shi2015extra}
		W.~Shi, Q.~Ling, G.~Wu, and W.~Yin.
		\newblock {EXTRA}: an exact first-order algorithm for decentralized consensus optimization.
		\newblock \emph{SIAM Journal on Optimization}, 25(2):944--966, 2015.
		
		\bibitem[Yuan et~al.(2019)Yuan, Ying, and Sayed]{yuan2019exactdiffusion}
		K.~Yuan, B.~Ying, and A.~H. Sayed.
		\newblock Exact diffusion for distributed optimization and learning---Part {I}: algorithm development.
		\newblock \emph{IEEE Transactions on Signal Processing}, 67(3):708--723, 2019.
		
	\end{thebibliography}
	
	%%%%%%%%%%%%%%%%%%%%%%%%%%%%%%%%%%%%%%%%%%%%%%%%%%%%%%%%%%%%
	
	\appendix
	\section{Dispersion Recursions, Burn-in Enforcement, and Logarithmic Mixing}\label{app:dispersion}
	
	This appendix provides a complete enforcement argument for the explicit logarithmic schedule~\eqref{eq:log_schedule}.
	The goal is to bound the network-induced error terms appearing in
	Proposition~\ref{prop:inexact-rn}, Lemma~\ref{lem:lambda-comparable},
	by quantities that decay faster than the target stationarity level $\varepsilon$.
	Throughout this appendix we fix $p>2$ and set
	\[
	\tau_k=t_k=\left\lceil\frac{p\log(k+2)}{-\log\rho}\right\rceil,\qquad k\ge 0,
	\]
	so that
	\begin{equation}\label{eq:rho_poly}
		\rho^{\tau_k}\le (k+2)^{-p},\qquad \rho^{t_k}\le (k+2)^{-p},\qquad k\ge 0.
	\end{equation}
	
	\subsection{Auxiliary inequalities}
	
	\begin{lemma}[Mean minimizes squared deviations]\label{lem:mean_minimizer_app}
		Let $\{u_i\}_{i=1}^N\subset\R^m$ and define $\bar u:=\frac1N\sum_{i=1}^N u_i$.
		Then for any $v\in\R^m$,
		\[
		\frac1N\sum_{i=1}^N\norm{u_i-\bar u}^2\le \frac1N\sum_{i=1}^N\norm{u_i-v}^2.
		\]
	\end{lemma}
	\begin{proof}
		Expand the right-hand side:
		\begin{align*}
			\sum_{i=1}^N\norm{u_i-v}^2
			&=\sum_{i=1}^N\norm{(u_i-\bar u)+(\bar u-v)}^2\\
			&=\sum_{i=1}^N\Big(\norm{u_i-\bar u}^2+2\ip{u_i-\bar u}{\bar u-v}+\norm{\bar u-v}^2\Big)\\
			&=\sum_{i=1}^N\norm{u_i-\bar u}^2+2\ip{\sum_{i=1}^N(u_i-\bar u)}{\bar u-v}+N\norm{\bar u-v}^2.
		\end{align*}
		Since $\sum_{i=1}^N(u_i-\bar u)=0$, the cross term vanishes and we obtain
		\[
		\sum_{i=1}^N\norm{u_i-v}^2=\sum_{i=1}^N\norm{u_i-\bar u}^2+N\norm{\bar u-v}^2\ge \sum_{i=1}^N\norm{u_i-\bar u}^2.
		\]
		Divide by $N$.
	\end{proof}
	
	\begin{lemma}[Uniform bound on local steps]\label{lem:step_bound_app}
		Under Assumption~\ref{ass:standing}, define the (finite) constant
		\[
		G_g:=\sup_{k\ge 0}\ \max_{1\le i\le N}\ \norm{\tilde g_{i,k}}\ \in(0,\infty).
		\]
		(For numerical stability, $G_g$ can be enforced in practice by clipping $\tilde g_{i,k}$; it is used here only to simplify the dispersion recursion.)
		Then for every $k\ge 0$ and every $i$, the local step of Algorithm~\ref{alg:disgrem} satisfies
		\begin{equation}\label{eq:step_unif_app}
			\norm{s_{i,k}}\le \sqrt{\frac{\norm{\tilde g_{i,k}}}{M}}\le \sqrt{\frac{G_g}{M}}.
		\end{equation}
		Consequently, $\norm{S_k}\le \sqrt{NG_g/M}$ for all $k$.
	\end{lemma}
	\begin{proof}
		Fix $k$ and $i$. The local step satisfies
		\[
		(\tilde H_{i,k}+(\lambda_{i,k}+\hat e_{i,k})I)s_{i,k}=-\tilde g_{i,k},
		\qquad \lambda_{i,k}=\sqrt{M\norm{\tilde g_{i,k}}},\qquad \hat e_{i,k}\ge 0.
		\]
		By construction $\tilde H_{i,k}\succeq 0$ and hence
		$\tilde H_{i,k}+(\lambda_{i,k}+\hat e_{i,k})I\succeq (\lambda_{i,k}+\hat e_{i,k})I$.
		Therefore
		\[
		(\lambda_{i,k}+\hat e_{i,k})\norm{s_{i,k}}
		\le \norm{(\tilde H_{i,k}+(\lambda_{i,k}+\hat e_{i,k})I)s_{i,k}}=\norm{\tilde g_{i,k}}.
		\]
		Dropping $\hat e_{i,k}\ge 0$ yields $\lambda_{i,k}\norm{s_{i,k}}\le \norm{\tilde g_{i,k}}$.
		If $\tilde g_{i,k}=0$ then $s_{i,k}=0$ and the bound holds. Otherwise, divide by $\lambda_{i,k}>0$:
		\[
		\norm{s_{i,k}}\le \frac{\norm{\tilde g_{i,k}}}{\lambda_{i,k}}=\frac{\norm{\tilde g_{i,k}}}{\sqrt{M\norm{\tilde g_{i,k}}}}=\sqrt{\frac{\norm{\tilde g_{i,k}}}{M}}.
		\]
		The final inequality follows from the definition of $G_g$.
		Summing \eqref{eq:step_unif_app} over $i$ gives $\norm{S_k}^2=\sum_i\norm{s_{i,k}}^2\le N(G_g/M)$.
	\end{proof}
	
	\subsection{Dispersion measures and basic contractions}
	
	For a stacked vector $Z=[z_1;\dots;z_N]\in\R^{Nd}$, define its RMS dispersion
	\[
	D(Z):=\sqrt{\frac1N\sum_{i=1}^N\norm{z_i-\bar z}^2}=\frac{1}{\sqrt N}\norm{((\Pmat)\otimes I_d)Z}.
	\]
	Similarly, for stacked matrices $\mathcal H=[H_1;\dots;H_N]\in\R^{Nd^2}$ with Frobenius norm on each block,
	\[
	D_H(\mathcal H):=\sqrt{\frac1N\sum_{i=1}^N\normF{H_i-\bar H}^2}.
	\]
	
	\begin{lemma}[Contraction by mixing]\label{lem:mix_contract_app}
		Let $W$ satisfy Assumption~\ref{ass:standing}(3). Then for any stacked vector $Z\in\R^{Nd}$ and integer $t\ge 1$,
		\[
		D\big((W^t\otimes I_d)Z\big)\le \rho^t\,D(Z).
		\]
		Likewise, for any stacked matrices $\mathcal H\in\R^{Nd^2}$,
		\[
		D_H\big((W^t\otimes I_{d^2})\mathcal H\big)\le \rho^t\,D_H(\mathcal H).
		\]
	\end{lemma}
	\begin{proof}
		We prove the vector statement; the matrix statement is identical with $I_d$ replaced by $I_{d^2}$.
		Let $Z^\perp:=((\Pmat)\otimes I_d)Z$. Since $W$ is doubly stochastic, $\J W^t=\J$ and $W^t\J=\J$, hence
		$(\Pmat W^t)=(W^t-\J)$ and
		\[
		((\Pmat)\otimes I_d)(W^t\otimes I_d)Z=((W^t-\J)\otimes I_d)Z.
		\]
		\[
		\text{Note that }(W^t-\J)\J=0\text{ and }(W^t-\J)=(W^t-\J)\Pmat,  
		\]
		\[
		\text{ hence }((W^t-\J)\otimes I_d)Z=((W^t-\J)\otimes I_d)Z^\perp.
		\]
		By Lemma~\ref{lem:consensus}, $\norm{((W^t-\J)\otimes I_d)Z^\perp}\le \rho^t\norm{Z^\perp}$.
		Applying this to $Z^\perp$ and using $D(Z)=\norm{Z^\perp}/\sqrt N$ gives the claim.
	\end{proof}
	
	\subsection{Primal dispersion recursion and decay}
	
	Recall the stacked post-mixing update $X_{k+1}=(W^{t_k}\otimes I_d)(\tilde X_k+S_k)$ and the pre-mixing $\tilde X_{k+1}=(W^{\tau_{k+1}}\otimes I_d)X_{k+1}$.
	
	\begin{lemma}[Primal dispersion recursion]\label{lem:DX_rec_strict}
		For all $k\ge 0$,
		\begin{equation}\label{eq:DX_rec_strict}
			D(\tilde X_{k+1})\le \rho^{\tau_{k+1}}\rho^{t_k}\Big(D(\tilde X_k)+\frac{\norm{S_k}}{\sqrt N}\Big).
		\end{equation}
	\end{lemma}
	\begin{proof}
		First apply Lemma~\ref{lem:mix_contract_app} to the post-mixing step:
		\[
		D(X_{k+1})=D\big((W^{t_k}\otimes I_d)(\tilde X_k+S_k)\big)
		\le \rho^{t_k} D(\tilde X_k+S_k).
		\]
		By the triangle inequality and the definition of $D(\cdot)$,
		\[
		D(\tilde X_k+S_k)=\sqrt{\frac1N\sum_i\norm{(\tilde x_{i,k}-\bar x_k)+(s_{i,k}-\bar s_k)}^2}
		\le D(\tilde X_k)+\sqrt{\frac1N\sum_i\norm{s_{i,k}-\bar s_k}^2}.
		\]
		Since $\sum_i\norm{s_{i,k}-\bar s_k}^2\le \sum_i\norm{s_{i,k}}^2=\norm{S_k}^2$, we obtain
		$D(\tilde X_k+S_k)\le D(\tilde X_k)+\norm{S_k}/\sqrt N$.
		Thus
		\[
		D(X_{k+1})\le \rho^{t_k}\Big(D(\tilde X_k)+\frac{\norm{S_k}}{\sqrt N}\Big).
		\]
		Second apply Lemma~\ref{lem:mix_contract_app} to the pre-mixing step:
		\[
		D(\tilde X_{k+1})=D\big((W^{\tau_{k+1}}\otimes I_d)X_{k+1}\big)\le \rho^{\tau_{k+1}} D(X_{k+1}),
		\]
		and combine.
	\end{proof}
	
	\begin{lemma}[Post-mixing primal dispersion]\label{lem:DX_post}
		For all $k\ge 0$,
		\[
		D(X_{k+1})\le \rho^{t_k}\Big(D(\tilde X_k)+\frac{\|S_k\|}{\sqrt N}\Big).
		\]
	\end{lemma}
	\begin{proof}
		This is exactly the first inequality derived in the proof of Lemma~\ref{lem:DX_rec_strict}.
	\end{proof}
	
	\begin{lemma}[Post-mixing primal dispersion decay]\label{lem:DX_post_decay}
		There exists a constant $C_X^{+}>0$ such that for all $k\ge 0$,
		\[
		D(X_k)\le \frac{C_X^{+}}{(k+1)^{p}}.
		\]
	\end{lemma}
	\begin{proof}
		From Lemma~\ref{lem:DX_post}, $\rho^{t_k}\le (k+2)^{-p}$ under \eqref{eq:log_schedule}, and the uniform bounds
		$D(\tilde X_k)\le D_{\max}$ (as in Lemma~\ref{lem:DX_decay_strict}) and $\|S_k\|/\sqrt N\le B$ (Lemma~\ref{lem:step_bound_app}),
		we get $D(X_{k+1})\le (k+2)^{-p}(D_{\max}+B)$. Set $C_X^{+}:=D_{\max}+B$.
	\end{proof}
	
	\begin{lemma}[Primal dispersion decay]\label{lem:DX_decay_strict}
		There exists a constant $C_X>0$ such that for all $k\ge 0$,
		\begin{equation}\label{eq:DX_decay_strict}
			D(\tilde X_k)\le \frac{C_X}{(k+2)^{p-1}}.
		\end{equation}
	\end{lemma}
\begin{proof}
	Let $D_k:=D(\tilde X_k)$ and $B:=\sup_k\|S_k\|/\sqrt N\le \sqrt{G_g/M}$ from Lemma~\ref{lem:step_bound_app}.
	From Lemma~\ref{lem:DX_rec_strict} we have
	\[
	D_{k+1}\le a_k(D_k+B),\qquad a_k:=\rho^{\tau_{k+1}}\rho^{t_k}.
	\]
	Under the schedule \eqref{eq:log_schedule}, $t_k\ge \left\lceil \frac{p\log 2}{-\log\rho}\right\rceil$ for all $k$,
	hence $\rho^{t_k}\le 2^{-p}$. Therefore $a_k\le \rho^{t_k}\le q:=2^{-p}<1$ and
	\[
	D_{k+1}\le q D_k+qB.
	\]
	This implies the uniform bound $D_k\le q^kD_0+\frac{q}{1-q}B\le D_{\max}$ where
	$D_{\max}:=D_0+\frac{q}{1-q}B$.
	Returning to $a_k\le (k+2)^{-p}$ from \eqref{eq:rho_poly}, we obtain
	\[
	D_{k+1}\le (k+2)^{-p}(D_{\max}+B)\le \frac{D_{\max}+B}{(k+2)^{p-1}}.
	\]
	Setting $C_X:=D_{\max}+B$ yields \eqref{eq:DX_decay_strict}.
\end{proof}
	
	
	\subsection{Gradient-tracker dispersion recursion and decay}
	
	We derive a recursion for the pre-mixed gradient-tracker dispersion $D(\tilde G_k)$.
	
	\begin{lemma}[Pointwise gradient disagreement bound]\label{lem:grad_disagree_pointwise_app}
		Let $Z=[z_1;\dots;z_N]\in\R^{Nd}$ and $\bar z:=\frac1N\sum_{i=1}^N z_i$.
		Then
		\begin{equation}\label{eq:grad_disagree_pointwise_app}
			\frac{1}{\sqrt N}\norm{((\Pmat)\otimes I_d)\,\nabla F(Z)}
			\le L_1\,D(Z)+\frac{\sigma_g}{\sqrt N}.
		\end{equation}
	\end{lemma}
	\begin{proof}
		Add and subtract $\nabla F(\one\otimes\bar z)$ and use the triangle inequality:
		\[
		((\Pmat)\otimes I_d)\nabla F(Z)=((\Pmat)\otimes I_d)\big(\nabla F(Z)-\nabla F(\one\otimes\bar z)\big)+((\Pmat)\otimes I_d)\nabla F(\one\otimes\bar z).
		\]
		For the second term, by definition of $\sigma_g$ in \eqref{eq:sigma_g_def},
		\[
		\norm{((\Pmat)\otimes I_d)\nabla F(\one\otimes\bar z)}\le \sigma_g.
		\]
		For the first term, use $\norm{((\Pmat)\otimes I_d)u}\le \norm{u}$ and $L_1$-Lipschitzness:
		\[
		\norm{\nabla F(Z)-\nabla F(\one\otimes\bar z)}^2
		=\sum_{i=1}^N\norm{\nabla f_i(z_i)-\nabla f_i(\bar z)}^2
		\le \sum_{i=1}^N (L_1\norm{z_i-\bar z})^2
		=L_1^2\norm{Z^\perp}^2,
		\]
		where $Z^\perp:=((\Pmat)\otimes I_d)Z$.
		Taking square roots and dividing by $\sqrt N$ yields \eqref{eq:grad_disagree_pointwise_app}.
	\end{proof}
	
	\begin{lemma}[Gradient-tracker dispersion recursion]\label{lem:DG_rec_strict}
		For all $k\ge 0$,
		\begin{equation}\label{eq:DG_rec_strict}
D(\tilde G_{k+1})
\le \rho^{\tau_{k+1}}\rho^{t_k}\Big(D(\tilde G_k)+L_1 D(X_{k+1})+L_1 D(X_k)+\frac{2\sigma_g}{\sqrt N}\Big).
		\end{equation}
	\end{lemma}
	\begin{proof}
		Write the stacked update for the gradient tracker (Algorithm~\ref{alg:disgrem}):
		\[
G_{k+1}=(W^{t_k}\otimes I_d)\Big(\tilde G_k+\nabla F(X_{k+1})-\nabla F(X_k)\Big),
\qquad
\tilde G_{k+1}=(W^{\tau_{k+1}}\otimes I_d)G_{k+1}.
		\]
		Apply Lemma~\ref{lem:mix_contract_app} to the pre-mixing: $D(\tilde G_{k+1})\le \rho^{\tau_{k+1}}D(G_{k+1})$.
		For $D(G_{k+1})$, apply $(\Pmat\otimes I_d)$, use $(\Pmat W^{t_k})=(W^{t_k}-\J)$, and take norms to get
		\[
D(G_{k+1})\le \rho^{t_k}\Big(D(\tilde G_k)+\frac{1}{\sqrt N}\norm{((\Pmat)\otimes I_d)(\nabla F(X_{k+1})-\nabla F(X_k))}\Big).
		\]
		Bound the forcing term by the triangle inequality and Lemma~\ref{lem:grad_disagree_pointwise_app} applied to $Z=X_{k+1}$ and $Z=X_k$.
		Substitute back to obtain \eqref{eq:DG_rec_strict}.
	\end{proof}
	
	\begin{lemma}[Gradient-tracker dispersion decay]\label{lem:DG_decay_strict}
		There exists a constant $C_G>0$ such that for all $k\ge 0$,
		\begin{equation}\label{eq:DG_decay_strict}
			D(\tilde G_k)\le \frac{C_G}{(k+2)^{p-1}}.
		\end{equation}
	\end{lemma}
\begin{proof}
	Let $G_k:=D(\tilde G_k)$. From Lemma~\ref{lem:DG_rec_strict} and \eqref{eq:rho_poly},
	\[
	G_{k+1}\le a_k\big(G_k+B_G\big),\qquad a_k:=\rho^{\tau_{k+1}}\rho^{t_k}\le q:=2^{-p}<1,
	\]
	where under \eqref{eq:bounded_levelset} we may take $B_G:=4L_1D+2\sigma_g/\sqrt N$.
	Thus $G_{k+1}\le qG_k+qB_G$, implying $G_k\le G_{\max}$ with
	$G_{\max}:=G_0+\frac{q}{1-q}B_G$.
	Using again $a_k\le (k+2)^{-p}$ yields
	\[
	G_{k+1}\le (k+2)^{-p}(G_{\max}+B_G)\le \frac{G_{\max}+B_G}{(k+2)^{p-1}}.
	\]
	Set $C_G:=G_{\max}+B_G$.
\end{proof}
	
	\subsection{Hessian-tracker dispersion recursion and decay}
	
	We now control the pre-mixed Hessian-tracker dispersion $D_H(\tilde H_k)$.
	
	\begin{lemma}[Pointwise Hessian disagreement bound]\label{lem:hess_disagree_pointwise_app}
		Let $Z=[z_1;\dots;z_N]\in\R^{Nd}$ and $\bar z:=\frac1N\sum_{i=1}^N z_i$.
		Then
		\begin{equation}\label{eq:hess_disagree_pointwise_app}
			\frac{1}{\sqrt N}\normF{((\Pmat)\otimes I_{d^2})\,\nabla^2 F(Z)}
			\le \sqrt d\,L_2\,D(Z)+\sigma_H.
		\end{equation}
	\end{lemma}
	\begin{proof}
		Add and subtract $\nabla^2 F(\one\otimes\bar z)$ and apply the triangle inequality:
		\[
		\normF{((\Pmat)\otimes I_{d^2})\,\nabla^2 F(Z)}
		\le \normF{((\Pmat)\otimes I_{d^2})\big(\nabla^2 F(Z)-\nabla^2 F(\one\otimes\bar z)\big)}
		+\normF{((\Pmat)\otimes I_{d^2})\,\nabla^2 F(\one\otimes\bar z)}.
		\]
		The second term is bounded by $\sqrt N\,\sigma_H$ by the definition \eqref{eq:sigma_H_def} (since $\norm{\bar z-x_\star}\le D$ whenever each $z_i$ satisfies \eqref{eq:bounded_levelset}).
		For the first term, use $\normF{A}\le \sqrt d\,\normtwo{A}$ and $L_2$-Lipschitzness of $\nabla^2 f_i$:
		\[
		\normF{\nabla^2 F(Z)-\nabla^2 F(\one\otimes\bar z)}^2
		=\sum_{i=1}^N\normF{\nabla^2 f_i(z_i)-\nabla^2 f_i(\bar z)}^2
		\le \sum_{i=1}^N\big(\sqrt d\,L_2\norm{z_i-\bar z}\big)^2
		=dL_2^2\norm{Z^\perp}^2,
		\]
		where $Z^\perp:=((\Pmat)\otimes I_d)Z$. Divide by $\sqrt N$ and use $\norm{Z^\perp}=\sqrt N\,D(Z)$.
	\end{proof}
	
	\begin{lemma}[Hessian-tracker dispersion recursion]\label{lem:DH_rec_strict}
		For all $k\ge 0$,
		\begin{equation}\label{eq:DH_rec_strict}
D_H(\tilde H_{k+1})
\le \rho^{\tau_{k+1}}\rho^{t_k}\Big(D_H(\tilde H_k)+\sqrt d\,L_2 D(X_{k+1})+\sqrt d\,L_2 D(X_k)+2\sigma_H\Big).
		\end{equation}
	\end{lemma}
	\begin{proof}
		Use the stacked Hessian tracker update, the pre-mixing contraction, and the non-expansiveness of PSD projection (Lemma~\ref{lem:psd-dispersion}) as in Lemma~\ref{lem:DG_rec_strict}.
		The forcing term is bounded by applying Lemma~\ref{lem:hess_disagree_pointwise_app} to $Z=X_{k+1}$ and $Z=X_k$.
	\end{proof}
	
	\begin{lemma}[Hessian-tracker dispersion decay]\label{lem:DH_decay_strict}
		There exists a constant $C_H>0$ such that for all $k\ge 0$,
		\begin{equation}\label{eq:DH_decay_strict}
			D_H(\tilde H_k)\le \frac{C_H}{(k+2)^{p-1}}.
		\end{equation}
	\end{lemma}
\begin{proof}
	Let $H_k:=D_H(\tilde H_k)$. From Lemma~\ref{lem:DH_rec_strict} and \eqref{eq:rho_poly},
	\[
	H_{k+1}\le a_k\big(H_k+B_H\big),\qquad a_k:=\rho^{\tau_{k+1}}\rho^{t_k}\le q:=2^{-p}<1,
	\]
	where using \eqref{eq:bounded_levelset} we may take $B_H:=4\sqrt d\,L_2D+2\sigma_H$.
	Thus $H_{k+1}\le qH_k+qB_H$, implying $H_k\le H_{\max}$ with
	$H_{\max}:=H_0+\frac{q}{1-q}B_H$.
	Using again $a_k\le (k+2)^{-p}$ yields
	\[
	H_{k+1}\le (k+2)^{-p}(H_{\max}+B_H)\le \frac{H_{\max}+B_H}{(k+2)^{p-1}}.
	\]
	Set $C_H:=H_{\max}+B_H$.
\end{proof}
	\begin{proof}[Proof of Proposition~\ref{prop:dispersion_decay}]
		Combine Lemmas~\ref{lem:DX_decay_strict}, \ref{lem:DG_decay_strict}, and \ref{lem:DH_decay_strict}.
		The final statement follows from Cauchy--Schwarz together with $\Delta_k^g\le D(\tilde G_k)$ and $\Delta_k^H\le D_H(\tilde H_k)$.
	\end{proof}
	
	\subsection{Burn-in index for enforcing relative conditions}
	
\begin{lemma}[A concrete burn-in index (hard-constant version)]\label{lem:burnin_index}
	Let $C_X,C_G,C_H$ be the constants in Lemmas~\ref{lem:DX_decay_strict}, \ref{lem:DG_decay_strict}, and \ref{lem:DH_decay_strict},
	and let $C_X^{+}$ be the constant in Lemma~\ref{lem:DX_post_decay}.
	Fix $\varepsilon>0$ and set $\eta:=1/12$ and $\alpha_b=\alpha_d:=1/4$.
	
	Define the hard lower-bound proxy
	\[
	\lambda_\varepsilon := \sqrt{\frac{3M\varepsilon}{5}},
	\qquad
	A_\varepsilon := M_H + \Big(1+\frac{\eta}{8}\Big)\lambda_\varepsilon
	= M_H + \frac{97}{96}\lambda_\varepsilon.
	\]
	Define two hard tolerances
	\[
	T_1(\varepsilon)
	:= \frac{\varepsilon\,\lambda_\varepsilon}{240\,(M_H+\lambda_\varepsilon)\,A_\varepsilon},
	\qquad
	T_2(\varepsilon)
	:= \frac{\varepsilon\,\lambda_\varepsilon}{240\,A_\varepsilon}.
	\]
	
	Define the thresholds
	\[
	\delta_X(\varepsilon)
	:= \min\Big\{\frac{\varepsilon}{64L_1},\ \frac{\eta}{16}\sqrt{\frac{3}{5}}\cdot\frac{\sqrt{M\varepsilon}}{L_2},\ \frac{T_2(\varepsilon)}{8L_1}\Big\},
	\qquad
\delta_X^{+}(\varepsilon)
:= \min\Big\{\frac{\varepsilon}{64L_1},\ \frac{T_2(\varepsilon)}{4L_1}\Big\},
\qquad
\delta_G(\varepsilon)
:= \min\Big\{
\frac{T_2(\varepsilon)}{4},\
\frac{3M\,\varepsilon^2}{8400\,(M_H+\lambda_\varepsilon)\,A_\varepsilon},\
\frac{\varepsilon}{64\sqrt N}
\Big\},
	\qquad
	\delta_H(\varepsilon)
	:= \frac{M\,T_1(\varepsilon)}{4}.
	\]
	
	Define the burn-in index
	\begin{equation}\label{eq:K0_def}
		K_0(\varepsilon)
		:=\min\Big\{k\ge 0:\ 
		\frac{C_X}{(k+2)^{p-1}}\le \delta_X(\varepsilon),\ 
		\frac{C_X^{+}}{(k+1)^{p}}\le \delta_X^{+}(\varepsilon),\ 
		\frac{C_G}{(k+2)^{p-1}}\le \delta_G(\varepsilon),\ 
		\frac{C_H}{(k+2)^{p-1}}\le \delta_H(\varepsilon)
		\Big\}.
	\end{equation}
	Then for all $k\ge K_0(\varepsilon)$,
	\[
	D(\tilde X_k)\le \delta_X(\varepsilon),\qquad
	D(X_k)\le \delta_X^{+}(\varepsilon),\qquad
	D(\tilde G_k)\le \delta_G(\varepsilon),\qquad
	D_H(\tilde H_k)\le \delta_H(\varepsilon).
	\]
\end{lemma}
	\begin{proof}
		Immediate from the decay bounds and the definition of $K_0(\varepsilon)$.
	\end{proof}
	\noindent\textbf{How this is used.}
	On any iteration $k$ with $\norm{\nabla f(\bar x_k)}\ge \varepsilon$,
	the two components of $\delta_X(\varepsilon)$ ensure simultaneously:
	(i) the \emph{gradient-bridge} forcing terms are $O(\varepsilon)$, and
	(ii) the \emph{Hessian-bridge} forcing term is $O(\sqrt{\varepsilon})$, matching the scale of $\lambda_k\asymp\sqrt{M\|g_k\|}$.
	This resolves the $\varepsilon$ vs.\ $\sqrt{\varepsilon}$ mismatch that would otherwise break Proposition~\ref{prop:inexact-rn}(3).
	
	\begin{proof}[Proof of Proposition~\ref{prop:burnin_enforce}]
		Fix $p>2$ and use the logarithmic schedule~\eqref{eq:log_schedule}.
		Set $\eta:=1/12$ and $\alpha_b=\alpha_d=1/4$.
		Let $\varepsilon\ge \varepsilon_{\mathrm{floor}}$ and let $K_0(\varepsilon)$ be as in Lemma~\ref{lem:burnin_index}.
		Fix any $k\ge K_0(\varepsilon)$ such that $\norm{\nabla f(\bar x_k)}\ge \varepsilon$.
		
		\medskip
		\noindent\textbf{1) Relative bridge.}
		By Lemma~\ref{lem:bridge-grad} and Lemma~\ref{lem:tildeDelta-upper},
		\[
		\norm{g_k-\tilde{\bar g}_k}
		\le D(\tilde G_k)+2L_1D(\tilde X_k)+L_1D(X_k)+\frac{\sigma_g}{\sqrt N}.
		\]
		Since $k\ge K_0(\varepsilon)$, Lemma~\ref{lem:burnin_index} yields $D(\tilde G_k)\le \varepsilon/(64\sqrt N)$ and $D(\tilde X_k)\le \varepsilon/(64L_1)$. Moreover, Lemma~\ref{lem:burnin_index} also gives
		\[
		D(X_k)\le \delta_X^{+}(\varepsilon)\le \frac{\varepsilon}{64L_1}.
		\]
		Also $\varepsilon\ge 16\sigma_g/\sqrt N$ gives $\sigma_g/\sqrt N\le \varepsilon/16$.
		
		Therefore,
		\[
		\norm{g_k-\tilde{\bar g}_k}
		\le \frac{\varepsilon}{64\sqrt N}+\frac{\varepsilon}{32}+\frac{\varepsilon}{64}+\frac{\varepsilon}{16}
		\le \frac{\varepsilon}{8},
		\]
		and hence $\norm{\tilde{\bar g}_k}\ge 7\varepsilon/8$. Thus \eqref{eq:rel-bridge} holds with $\alpha_b=1/4$.
		
		\medskip
		\noindent\textbf{2) Relative dispersion.}
		Using $\max_i\norm{\tilde g_{i,k}-\tilde{\bar g}_k}\le \sqrt N D(\tilde G_k)$ and $D(\tilde G_k)\le \varepsilon/(64\sqrt N)$ gives
		$\max_i\norm{\tilde g_{i,k}-\tilde{\bar g}_k}\le \varepsilon/64\le (1/56)\norm{\tilde{\bar g}_k}$.
		Hence \eqref{eq:rel-disp} holds with $\alpha_d=1/4$.
		
		\medskip
		\noindent\textbf{3) Hessian absorption and bridge accuracy.}
		Let $\bar\lambda_k:=\frac1N\sum_i\sqrt{M\norm{\tilde g_{i,k}}}$.
		By \eqref{eq:rel-disp}, $\norm{\tilde g_{i,k}}\ge (1-\alpha_d)\norm{\tilde{\bar g}_k}$ for all $i$, so
		$\bar\lambda_k\ge \sqrt{M(1-\alpha_d)\norm{\tilde{\bar g}_k}}$.
		Using \eqref{eq:rel-bridge} yields $\norm{\tilde{\bar g}_k}\ge \norm{g_k}/(1+\alpha_b)$ and thus
		\[
		\bar\lambda_k\ge \sqrt{\frac{3}{5}}\sqrt{M\norm{g_k}}\ge \sqrt{\frac{3}{5}}\sqrt{M\varepsilon}.
		\]
		Since $\varepsilon\ge 61440\,\bar e_{\max}^2/M$, we have $\bar\lambda_k\ge 192\,\bar e_{\max}\ge 192\,\bar e_k$.
		Let $\alpha:=\eta/16=1/192$. Then $191\,\bar e_k\le \bar\lambda_k$ implies $(1-\alpha)\bar e_k\le \alpha\bar\lambda_k$, which is equivalent to
		\[
		\bar e_k\le \alpha(\bar\lambda_k+\bar e_k)=\frac{\eta}{16}\lambda_k
		\quad\Longrightarrow\quad
		\bar e_k\le \frac{\eta}{8}\lambda_k,
		\]
		i.e., \eqref{eq:e-absorb} holds (in fact with a stronger constant).
		
		Finally, Lemma~\ref{lem:bridge-hess} gives
		$\normtwo{\nabla^2 f(\bar x_k)-\tilde{\bar H}_k}\le L_2D(\tilde X_k)+\bar e_k$.
		By Lemma~\ref{lem:burnin_index}, $D(\tilde X_k)\le \frac{\eta}{16}\sqrt{\frac{3}{5}}\cdot\frac{\sqrt{M\varepsilon}}{L_2}$, hence
		\[
		L_2D(\tilde X_k)\le \frac{\eta}{16}\sqrt{\frac{3}{5}}\sqrt{M\varepsilon}\le \frac{\eta}{16}\bar\lambda_k\le \frac{\eta}{16}\lambda_k.
		\]
		Combining with $\bar e_k\le (\eta/16)\lambda_k$ yields
		$L_2D(\tilde X_k)+\bar e_k\le (\eta/8)\lambda_k$, i.e., Proposition~\ref{prop:inexact-rn}(3).
		
		\medskip
	\medskip
\medskip
\noindent\textbf{4) Remaining inexact RN conditions (items (1) and (2) of Proposition~\ref{prop:inexact-rn}).}
Fix $k\ge K_0(\varepsilon)$ such that $\|\nabla f(\bar x_k)\|\ge \varepsilon$.
We prove items (1) and (2) of Proposition~\ref{prop:inexact-rn} with $\eta=1/12$ and $\alpha_b=\alpha_d=1/4$.

\smallskip
\noindent\underline{\emph{Step 4.1: hard lower bounds on $\underline\lambda_k$ and on the reference-system size.}}
By items (1)--(2) proved earlier in this proposition, the relative bridge and dispersion conditions
\eqref{eq:rel-bridge}--\eqref{eq:rel-disp} hold with $\alpha_b=\alpha_d=1/4$.
Hence
\[
\|\tilde{\bar g}_k\|\ge \frac{1}{1+\alpha_b}\|g_k\|\ge \frac{4}{5}\varepsilon,
\qquad
\min_i\|\tilde g_{i,k}\|\ge (1-\alpha_d)\|\tilde{\bar g}_k\|\ge \frac{3}{5}\varepsilon.
\]
Therefore, for every $i$,
\[
\lambda_{i,k}=\sqrt{M\|\tilde g_{i,k}\|}\ge \sqrt{\frac{3M\varepsilon}{5}}=\lambda_\varepsilon,
\qquad
\tilde\lambda_{i,k}=\lambda_{i,k}+\hat e_{i,k}\ge \lambda_{i,k}\ge \lambda_\varepsilon,
\]
and thus
\begin{equation}\label{eq:hard_lambda_lb_step4}
	\underline\lambda_k:=\min_i\tilde\lambda_{i,k}\ge \lambda_\varepsilon,
	\qquad
	\lambda_k:=\tilde{\bar\lambda}_k=\frac1N\sum_i\tilde\lambda_{i,k}\ge \lambda_\varepsilon.
\end{equation}

Moreover, by Proposition~\ref{prop:inexact-rn}(3) already established in Step~3,
\[
\|\nabla^2 f(\bar x_k)-\tilde{\bar H}_k\|\le \frac{\eta}{8}\lambda_k.
\]
Since $\|\nabla^2 f(\bar x_k)\|\le M_H$, we obtain
\[
\|\tilde{\bar H}_k\|\le M_H+\frac{\eta}{8}\lambda_k,
\qquad
\|A_k^{\mathrm{ref}}\|
=\|\tilde{\bar H}_k+\lambda_k I\|
\le \|\tilde{\bar H}_k\|+\lambda_k
\le M_H+\Big(1+\frac{\eta}{8}\Big)\lambda_k.
\]
Hence, using $\lambda_k\ge \lambda_\varepsilon$ and $\|\tilde{\bar g}_k\|\ge \frac45\varepsilon$,
\begin{equation}\label{eq:sref_hard_lb_step4}
	\|s_k^{\mathrm{ref}}\|
	=\|(A_k^{\mathrm{ref}})^{-1}\tilde{\bar g}_k\|
	\ge \frac{\|\tilde{\bar g}_k\|}{\|A_k^{\mathrm{ref}}\|}
	\ge \frac{\frac45\varepsilon}{M_H+\big(1+\frac{\eta}{8}\big)\lambda_k}
	\ge \frac{\frac45\varepsilon}{M_H+\big(1+\frac{\eta}{8}\big)\lambda_\varepsilon}
	=\frac{\frac45\varepsilon}{A_\varepsilon}.
\end{equation}

\smallskip
\noindent\underline{\emph{Step 4.2: item (1) (dispersion control).}}
By Lemma~\ref{lem:step-disp},
\[
\|\bar s_k-s_k^{\mathrm{ref}}\|
\le
\frac{1}{\underline\lambda_k}\Delta_k^g
+\frac{\|\tilde{\bar g}_k\|}{\underline\lambda_k\,\lambda_k}\big(\Delta_k^H+\Delta_k^\lambda\big).
\]
Using $\lambda_k\ge \underline\lambda_k$ and $\underline\lambda_k^2\ge M(1-\alpha_d)\|\tilde{\bar g}_k\|$ (from $\tilde\lambda_{i,k}\ge \lambda_{i,k}$ and \eqref{eq:rel-disp}),
we get the hard bound
\[
\frac{\|\tilde{\bar g}_k\|}{\underline\lambda_k\,\lambda_k}
\le
\frac{\|\tilde{\bar g}_k\|}{\underline\lambda_k^2}
\le \frac{1}{M(1-\alpha_d)}=\frac{4}{3M}.
\]
Moreover, Lemma~\ref{lem:delta_lambda_control} with $\alpha_d=1/4$ gives
\[
\Delta_k^\lambda
\le
\frac{\sqrt{M}}{\sqrt{1-\alpha_d}}\cdot\frac{\Delta_k^g}{\sqrt{\|\tilde{\bar g}_k\|}}
+2\bar e_k
\le
\sqrt{\frac{5M}{3}}\cdot\frac{\Delta_k^g}{\sqrt{\varepsilon}}+2\bar e_k,
\]
where we used $\|\tilde{\bar g}_k\|\ge \frac45\varepsilon$.
Combining the above and using $\underline\lambda_k\ge \lambda_\varepsilon$ yields
\begin{align*}
	\|\bar s_k-s_k^{\mathrm{ref}}\|
	&\le
	\frac{1}{\lambda_\varepsilon}\Delta_k^g
	+\frac{4}{3M}\Delta_k^H
	+\frac{4}{3M}\Delta_k^\lambda\\
	&\le
	\Big(\frac{1}{\lambda_\varepsilon}+\frac{4}{3M}\sqrt{\frac{5M}{3}}\cdot\frac{1}{\sqrt{\varepsilon}}\Big)\Delta_k^g
	+\frac{4}{3M}\Delta_k^H
	+\frac{8}{3M}\bar e_k\\
	&=
	\frac{7}{3}\sqrt{\frac{5}{3}}\cdot\frac{\Delta_k^g}{\sqrt{M\varepsilon}}
	+\frac{4}{3M}\Delta_k^H
	+\frac{8}{3M}\bar e_k.
\end{align*}
By Cauchy--Schwarz, $\Delta_k^g\le D(\tilde G_k)$ and $\Delta_k^H\le D_H(\tilde H_k)$.
By Lemma~\ref{lem:burnin_index}, for $k\ge K_0(\varepsilon)$ we have
$D(\tilde G_k)\le \delta_G(\varepsilon)$ and $D_H(\tilde H_k)\le \delta_H(\varepsilon)$.
By the definitions of $\delta_G(\varepsilon)$ and $\delta_H(\varepsilon)$ in Lemma~\ref{lem:burnin_index}, and by $\bar e_k\le \bar e_{\max}$
together with $\varepsilon\ge \varepsilon_{\mathrm{floor}}$ from \eqref{eq:eps_floor},
the three terms above are each bounded by $T_1(\varepsilon)/3$.
Therefore,
\begin{equation}\label{eq:hard_step_disp_T1}
	\|\bar s_k-s_k^{\mathrm{ref}}\|\le T_1(\varepsilon).
\end{equation}
Using \eqref{eq:sref_hard_lb_step4} and $\lambda_k\ge \lambda_\varepsilon$,
\[
T_1(\varepsilon)
=\frac{\varepsilon\lambda_\varepsilon}{240(M_H+\lambda_\varepsilon)A_\varepsilon}
\le
\frac{\eta}{16}\cdot\frac{\lambda_k}{M_H+\lambda_k}\cdot \|s_k^{\mathrm{ref}}\|.
\]
Hence $\|\bar s_k-s_k^{\mathrm{ref}}\|\le \frac{\eta}{16}\frac{\lambda_k}{M_H+\lambda_k}\|s_k^{\mathrm{ref}}\|$.
Since $\frac{\eta}{16}<1/2$, we have $\|\bar s_k\|\ge (1-\eta/16)\|s_k^{\mathrm{ref}}\|$ and therefore
\[
\|\bar s_k-s_k^{\mathrm{ref}}\|
\le
\frac{\eta}{16}\frac{\lambda_k}{M_H+\lambda_k}\cdot\frac{1}{1-\eta/16}\|\bar s_k\|
\le
\frac{\eta}{8}\cdot\frac{\lambda_k}{M_H+\lambda_k}\,\|\bar s_k\|,
\]
which is exactly item (1) of Proposition~\ref{prop:inexact-rn}.

\smallskip
\noindent\underline{\emph{Step 4.3: item (2) (gradient-bridge accuracy).}}
By Lemma~\ref{lem:tildeDelta-upper} and $\sqrt{\tilde E_k}=D(\tilde X_k)$,
\[
\frac{\|\tilde\Delta_k\|}{\sqrt N}+L_1\sqrt{\tilde E_k}
\le
D(\tilde G_k)+2L_1 D(\tilde X_k)+L_1 D(X_k)+\frac{\sigma_g}{\sqrt N}.
\]
By Lemma~\ref{lem:burnin_index}, for $k\ge K_0(\varepsilon)$ we have
\[
D(\tilde G_k)\le \delta_G(\varepsilon),\quad
D(\tilde X_k)\le \delta_X(\varepsilon),\quad
D(X_k)\le \delta_X^{+}(\varepsilon).
\]
By the definitions of $\delta_G(\varepsilon),\delta_X(\varepsilon),\delta_X^{+}(\varepsilon)$
and the condition $\varepsilon\ge \varepsilon_{\mathrm{floor}}$ in \eqref{eq:eps_floor},
each of the four terms above is bounded by $T_2(\varepsilon)/4$, hence
\begin{equation}\label{eq:hard_bridge_T2}
	\frac{\|\tilde\Delta_k\|}{\sqrt N}+L_1\sqrt{\tilde E_k}\le T_2(\varepsilon).
\end{equation}
Finally, \eqref{eq:hard_step_disp_T1} implies $\|\bar s_k\|\ge (1-\eta/16)\|s_k^{\mathrm{ref}}\|$,
so using $\lambda_k\ge \lambda_\varepsilon$ and \eqref{eq:sref_hard_lb_step4} we get
\[
\frac{\eta}{8}\lambda_k\|\bar s_k\|
\ge
\frac{\eta}{8}(1-\eta/16)\lambda_\varepsilon\|s_k^{\mathrm{ref}}\|
\ge
\frac{\eta}{8}(1-\eta/16)\lambda_\varepsilon\cdot \frac{\frac45\varepsilon}{A_\varepsilon}
\ge
\frac{\eta}{16}\lambda_\varepsilon\cdot \frac{\frac45\varepsilon}{A_\varepsilon}
=
T_2(\varepsilon),
\]
where we used $(1-\eta/16)\ge 1/2$.
Combining with \eqref{eq:hard_bridge_T2} yields
\[
\frac{\|\tilde\Delta_k\|}{\sqrt N}+L_1\sqrt{\tilde E_k}\le \frac{\eta}{8}\lambda_k\|\bar s_k\|,
\]
which is exactly item (2) of Proposition~\ref{prop:inexact-rn}.
This completes Step~4.
	\end{proof}
	
	
	\section{Proofs}\label{app:proofs}
	
	
	\begin{proof}[Proof of Lemma~\ref{lem:taylor3}]
		Fix $x,s$ and define $\phi(t):=h(x+ts)$ for $t\in[0,1]$.
		Then $\phi'(t)=\ip{\nabla h(x+ts)}{s}$ and $\phi''(t)=s^\top\nabla^2 h(x+ts)s$.
		By $L_2$-Lipschitzness of $\nabla^2 h$,
		\[
		|\phi''(t)-\phi''(0)|
		=|s^\top(\nabla^2 h(x+ts)-\nabla^2 h(x))s|
		\le \normtwo{\nabla^2 h(x+ts)-\nabla^2 h(x)}\,\norm{s}^2
		\le L_2 t\norm{s}^3.
		\]
		Hence $\phi''(t)\le \phi''(0)+L_2 t\norm{s}^3$. Integrate:
		\[
		\phi(1)=\phi(0)+\phi'(0)+\int_0^1 (1-t)\phi''(t)\,dt
		\le \phi(0)+\phi'(0)+\int_0^1(1-t)\big(\phi''(0)+L_2 t\norm{s}^3\big)\,dt.
		\]
		Compute $\int_0^1(1-t)\,dt=\frac12$ and $\int_0^1(1-t)t\,dt=\frac16$, giving
		\[
		\phi(1)\le \phi(0)+\phi'(0)+\tfrac12\phi''(0)+\tfrac{L_2}{6}\norm{s}^3.
		\]
		Substitute $\phi(0)=h(x)$, $\phi'(0)=\ip{\nabla h(x)}{s}$, $\phi''(0)=s^\top\nabla^2 h(x)s$.
	\end{proof}
	
	\begin{proof}[Proof of Lemma~\ref{lem:consensus}]
		Since $W$ is symmetric, it is orthogonally diagonalizable and its eigenvalues lie in $[-1,1]$.
		Because $W$ is doubly stochastic, $\one$ is an eigenvector with eigenvalue $1$, and $\J$ is the orthogonal projector onto $\mathrm{span}\{\one\}$.
		Moreover $W\J=\J W=\J$.
		We prove $W^t-\J=(W-\J)^t$ by induction.
		For $t=1$ it is trivial.
		Assume $(W-\J)^t=W^t-\J$. Then
		\[
		(W-\J)^{t+1}=(W-\J)^t(W-\J)=(W^t-\J)(W-\J)=W^{t+1}-W^t\J-\J W+\J^2.
		\]
		Using $W^t\J=\J$, $\J W=\J$, and $\J^2=\J$, we obtain $(W-\J)^{t+1}=W^{t+1}-\J$.
		Taking spectral norms:
		\[
		\normtwo{W^t-\J}=\normtwo{(W-\J)^t}\le \normtwo{W-\J}^t=\rho^t.
		\]
		For stacked vectors, $\normtwo{(A\otimes I_d)}=\normtwo{A}$.
	\end{proof}
	
	\begin{proof}[Proof of Lemma~\ref{lem:gap-grad}]
		By $L$-smoothness, for any $y$,
		\[
		h(y)\le h(x)+\ip{\nabla h(x)}{y-x}+\frac{L}{2}\norm{y-x}^2.
		\]
		Choose $y=x-\frac{1}{L}\nabla h(x)$:
		\[
		h\Big(x-\frac{1}{L}\nabla h(x)\Big)\le h(x)-\frac{1}{2L}\norm{\nabla h(x)}^2.
		\]
		Since $x_\star$ minimizes $h$, $h(x_\star)\le h(x-\frac{1}{L}\nabla h(x))$.
		Rearrange to get $\norm{\nabla h(x)}^2\le 2L(h(x)-h(x_\star))$.
	\end{proof}
	
	\begin{proof}[Proof of Lemma~\ref{lem:frob}]
		Let $\lambda_1,\dots,\lambda_d$ be eigenvalues of symmetric $E$.
		Then $\normtwo{E}=\max_j|\lambda_j|$ and $\normF{E}=\sqrt{\sum_j\lambda_j^2}$, hence $\normtwo{E}\le\normF{E}=\hat e$.
		Eigenvalues of $\hat e I\pm E$ are $\hat e\pm \lambda_j\ge \hat e-|\lambda_j|\ge 0$, so $\hat e I\pm E\succeq 0$.
		Thus $s^\top(\hat e I+E)s\ge 0$, i.e., $-\hat e\norm{s}^2-s^\top E s\le 0$.
	\end{proof}
	
	\begin{proof}[Proof of Lemma~\ref{lem:M-vs-s}]
		From the local linear system,
		\[
		(\tilde H_{i,k}+(\lambda_{i,k}+\hat e_{i,k})I)s_{i,k}=-\tilde g_{i,k}.
		\]
		Since $\tilde H_{i,k}\succeq 0$ and $\hat e_{i,k}\ge 0$, we have
		\[
		(\lambda_{i,k}+\hat e_{i,k})\norm{s_{i,k}} \le \norm{(\tilde H_{i,k}+(\lambda_{i,k}+\hat e_{i,k})I)s_{i,k}}
		= \norm{\tilde g_{i,k}},
		\]
		hence $\lambda_{i,k}\norm{s_{i,k}}\le \norm{\tilde g_{i,k}}$.
		Using $\lambda_{i,k}^2=M\norm{\tilde g_{i,k}}$, we get
		\[
		\lambda_{i,k}\norm{s_{i,k}}\le \frac{\lambda_{i,k}^2}{M}
		\quad\Longrightarrow\quad
		M\norm{s_{i,k}}\le \lambda_{i,k}.
		\]
		Since $M\ge L_2$, also $L_2\norm{s_{i,k}}\le \lambda_{i,k}$.
	\end{proof}
	
	\begin{proof}[Proof of Lemma~\ref{lem:local-descent}]
		Apply Lemma~\ref{lem:taylor3} to $h=f_i$ at $(\tilde x_{i,k},s_{i,k})$:
		\[
		f_i(\tilde x_{i,k}+s_{i,k})
		\le f_i(\tilde x_{i,k})+\ip{\nabla f_i(\tilde x_{i,k})}{s_{i,k}}
		+\frac12 s_{i,k}^\top\nabla^2 f_i(\tilde x_{i,k})s_{i,k}+\frac{L_2}{6}\norm{s_{i,k}}^3.
		\]
		Write $\nabla f_i(\tilde x_{i,k})=\tilde g_{i,k}-\tilde\delta_{i,k}$.
		Also from the step equation,
		\[
		\ip{\tilde g_{i,k}}{s_{i,k}}
		=-s_{i,k}^\top\tilde H_{i,k}s_{i,k}-(\lambda_{i,k}+\hat e_{i,k})\norm{s_{i,k}}^2.
		\]
		Substitute into the Taylor bound:
		\begin{align*}
			f_i(\tilde x_{i,k}+s_{i,k})
			&\le f_i(\tilde x_{i,k})
			+\ip{\tilde g_{i,k}}{s_{i,k}}-\ip{\tilde\delta_{i,k}}{s_{i,k}}
			+\frac12 s_{i,k}^\top\nabla^2 f_i(\tilde x_{i,k})s_{i,k}+\frac{L_2}{6}\norm{s_{i,k}}^3\\
			&= f_i(\tilde x_{i,k})
			-\ip{\tilde\delta_{i,k}}{s_{i,k}}
			-s_{i,k}^\top\tilde H_{i,k}s_{i,k}-(\lambda_{i,k}+\hat e_{i,k})\norm{s_{i,k}}^2
			+\frac12 s_{i,k}^\top\nabla^2 f_i(\tilde x_{i,k})s_{i,k}
			+\frac{L_2}{6}\norm{s_{i,k}}^3.
		\end{align*}
		Write $\tilde H_{i,k}=\nabla^2 f_i(\tilde x_{i,k})+\tilde e_{i,k}$:
		\[
		-s^\top\tilde H s + \frac12 s^\top\nabla^2 f_i s
		= -\frac12 s^\top\nabla^2 f_i s - s^\top \tilde e\, s \le - s^\top \tilde e\,s,
		\]
		since $\nabla^2 f_i\succeq 0$.
		By Lemma~\ref{lem:frob}, $-\hat e_{i,k}\norm{s}^2 - s^\top\tilde e\,s\le 0$, hence
		$- s^\top \tilde e\,s - \hat e_{i,k}\norm{s}^2 \le 0$.
		Dropping this nonpositive quantity yields
		\[
		f_i(\tilde x_{i,k}+s_{i,k})
		\le f_i(\tilde x_{i,k})
		-\lambda_{i,k}\norm{s_{i,k}}^2
		-\ip{\tilde\delta_{i,k}}{s_{i,k}}
		+\frac{L_2}{6}\norm{s_{i,k}}^3.
		\]
	\end{proof}
	
	\begin{proof}[Proof of Lemma~\ref{lem:local-clean}]
		From Lemma~\ref{lem:M-vs-s}, $L_2\norm{s_{i,k}}\le \lambda_{i,k}$, hence
		$\frac{L_2}{6}\norm{s_{i,k}}^3\le \frac{1}{6}\lambda_{i,k}\norm{s_{i,k}}^2$.
		Insert into \eqref{eq:local-descent} to obtain
		\[
		f_i(\tilde x_{i,k}+s_{i,k})
		\le f_i(\tilde x_{i,k})-\Big(1-\frac16\Big)\lambda_{i,k}\norm{s_{i,k}}^2-\ip{\tilde\delta_{i,k}}{s_{i,k}}.
		\]
	\end{proof}
	
	\begin{proof}[Proof of Lemma~\ref{lem:avg-pres}]
		Pre-mixing preserves averages since $W^{\tau_k}$ is doubly stochastic:
		\[
		\tilde{\bar x}_k=\frac1N\sum_i\tilde x_{i,k}
		=\frac1N\sum_i\sum_j [W^{\tau_k}]_{ij}x_{j,k}
		=\frac1N\sum_j \Big(\sum_i [W^{\tau_k}]_{ij}\Big)x_{j,k}
		=\frac1N\sum_j x_{j,k}=\bar x_k.
		\]
		Similarly, post-mixing is doubly stochastic and $y_{i,k+1}=\tilde x_{i,k}+s_{i,k}$, so
		\[
		\bar x_{k+1}=\frac1N\sum_i x_{i,k+1}
		=\frac1N\sum_i\sum_j[W^{t_k}]_{ij}y_{j,k+1}
		=\frac1N\sum_j y_{j,k+1}
		=\tilde{\bar x}_k+\bar s_k=\bar x_k+\bar s_k.
		\]
	\end{proof}
	
	\begin{proof}[Proof of Lemma~\ref{lem:avg-regime}]
		By Lemma~\ref{lem:M-vs-s}, $M\norm{s_{i,k}}\le \lambda_{i,k}$ for all $i$, hence
		\[
		\frac1N\sum_{i=1}^N \lambda_{i,k}\ge \frac{M}{N}\sum_{i=1}^N \norm{s_{i,k}}
		\ge M \norm{\frac1N\sum_{i=1}^N s_{i,k}}
		= M\norm{\bar s_k},
		\]
		using $\norm{\frac1N\sum_i s_i}\le \frac1N\sum_i\norm{s_i}$.
		Since $M\ge L_2$, $L_2\norm{\bar s_k}\le M\norm{\bar s_k}\le \frac1N\sum_i\lambda_{i,k}$.
		Finally, $\tilde{\bar\lambda}_k=\frac1N\sum_i(\lambda_{i,k}+\hat e_{i,k})\ge \frac1N\sum_i\lambda_{i,k}$.
	\end{proof}
	
	
	\begin{proof}[Proof of Lemma~\ref{lem:tildeDelta-upper}]
		Recall $\tilde\Delta_k=\tilde G_k-\nabla F(\tilde X_k)$ and $\bar x_k=\frac1N\sum_i \tilde x_{i,k}$.
		Add and subtract $\one\otimes\tilde{\bar g}_k$ and $\nabla F(\one\otimes\bar x_k)$:
		\[
		\tilde\Delta_k
		=\big(\tilde G_k-(\one\otimes\tilde{\bar g}_k)\big)
		+\big((\one\otimes\tilde{\bar g}_k)-\nabla F(\one\otimes\bar x_k)\big)
		+\big(\nabla F(\one\otimes\bar x_k)-\nabla F(\tilde X_k)\big).
		\]
		Divide by $\sqrt N$ and bound each term.
		
		\noindent\textbf{(i) Dispersion term.}
		By definition,
		$\frac{1}{\sqrt N}\|\tilde G_k-(\one\otimes\tilde{\bar g}_k)\|=D(\tilde G_k)$.
		
		\noindent\textbf{(ii) Heterogeneity + disagreement term.}
		Using Lemma~\ref{lem:avg-track}, $\tilde{\bar g}_k=\frac1N\sum_{i=1}^N\nabla f_i(x_{i,k})$.
		Hence
		\[
		\frac{1}{\sqrt N}\|(\one\otimes\tilde{\bar g}_k)-\nabla F(\one\otimes\bar x_k)\|
		\le \|\tilde{\bar g}_k-\nabla f(\bar x_k)\|+\frac{1}{\sqrt N}\|((\Pmat)\otimes I_d)\nabla F(\one\otimes\bar x_k)\|.
		\]
		The second term is bounded by $\sigma_g/\sqrt N$ by \eqref{eq:sigma_g_def}.
		For the first term,
		\[
		\|\tilde{\bar g}_k-\nabla f(\bar x_k)\|
		=\Big\|\frac1N\sum_{i=1}^N\big(\nabla f_i(x_{i,k})-\nabla f_i(\bar x_k)\big)\Big\|
		\le \frac1N\sum_{i=1}^N L_1\|x_{i,k}-\bar x_k\|
		\le L_1 D_X(k).
		\]
		
		\noindent\textbf{(iii) Lipschitz term.}
		By block-wise $L_1$-Lipschitzness,
		\[
		\frac{1}{\sqrt N}\|\nabla F(\one\otimes\bar x_k)-\nabla F(\tilde X_k)\|
		\le L_1\sqrt{\tilde E_k}.
		\]
		Combine (i)--(iii). Finally, $\Delta_k^g\le D(\tilde G_k)$ follows from Cauchy--Schwarz.
	\end{proof}
	
	
	\begin{proof}[Proof of Lemma~\ref{lem:bridge-grad}]
		Write
		\[
		\nabla f(\bar x_k)-\tilde{\bar g}_k
		=\frac1N\sum_{i=1}^N\big(\nabla f_i(\bar x_k)-\tilde g_{i,k}\big)
		=\frac1N\sum_{i=1}^N\big(\nabla f_i(\bar x_k)-\nabla f_i(\tilde x_{i,k})\big)
		-\frac1N\sum_{i=1}^N \tilde\delta_{i,k}.
		\]
		Take norms and apply triangle inequality:
		\[
		\norm{\nabla f(\bar x_k)-\tilde{\bar g}_k}
		\le \frac1N\sum_i\norm{\nabla f_i(\bar x_k)-\nabla f_i(\tilde x_{i,k})}
		+\frac1N\sum_i\norm{\tilde\delta_{i,k}}.
		\]
		By $L_1$-Lipschitzness,
		$\norm{\nabla f_i(\bar x_k)-\nabla f_i(\tilde x_{i,k})}\le L_1\norm{\bar x_k-\tilde x_{i,k}}$.
		By Cauchy--Schwarz,
		$\frac1N\sum_i\norm{\bar x_k-\tilde x_{i,k}}\le \sqrt{\frac1N\sum_i\norm{\bar x_k-\tilde x_{i,k}}^2}=\sqrt{\tilde E_k}$.
		Also,
		$\frac1N\sum_i\norm{\tilde\delta_{i,k}}\le \sqrt{\frac1N\sum_i\norm{\tilde\delta_{i,k}}^2}=\frac{\norm{\tilde\Delta_k}}{\sqrt N}$.
		Combine.
	\end{proof}
	
	\begin{proof}[Proof of Lemma~\ref{lem:bridge-hess}]
		Decompose
		\[
		\nabla^2 f(\bar x_k)-\tilde{\bar H}_k
		=\frac1N\sum_{i=1}^N\big(\nabla^2 f_i(\bar x_k)-\nabla^2 f_i(\tilde x_{i,k})\big)
		+\frac1N\sum_{i=1}^N\big(\nabla^2 f_i(\tilde x_{i,k})-\tilde H_{i,k}\big).
		\]
		Take operator norms and triangle inequality:
		\[
		\normtwo{\nabla^2 f(\bar x_k)-\tilde{\bar H}_k}
		\le \frac1N\sum_i\normtwo{\nabla^2 f_i(\bar x_k)-\nabla^2 f_i(\tilde x_{i,k})}
		+\frac1N\sum_i\normtwo{\nabla^2 f_i(\tilde x_{i,k})-\tilde H_{i,k}}.
		\]
		By $L_2$-Lipschitzness of Hessians,
		$\normtwo{\nabla^2 f_i(\bar x_k)-\nabla^2 f_i(\tilde x_{i,k})}\le L_2\norm{\bar x_k-\tilde x_{i,k}}$,
		and Cauchy--Schwarz gives the average bound $\le L_2\sqrt{\tilde E_k}$.
		For the second term, $\normtwo{A}\le \normF{A}$ implies
		$\normtwo{\nabla^2 f_i(\tilde x_{i,k})-\tilde H_{i,k}}\le \normF{\nabla^2 f_i(\tilde x_{i,k})-\tilde H_{i,k}}=\hat e_{i,k}$.
		Averaging yields $\bar e_k$.
	\end{proof}
	
	\begin{proof}[Proof of Remark~\ref{rem:psd-nonexp}]
		$\ProjPSD$ is the Euclidean projection onto a closed convex set in the Hilbert space $(\R^{d\times d},\ip{A}{B}=\mathrm{tr}(A^\top B))$,
		so it is firmly non-expansive; in particular, it is $1$-Lipschitz in Frobenius norm.
	\end{proof}
	
	\begin{proof}[Proof of Lemma~\ref{lem:psd-dispersion}]
		Let $T(A)=\ProjPSD(\mathrm{sym}(A))$. By Remark~\ref{rem:psd-nonexp}, $T$ is $1$-Lipschitz in Frobenius norm. Let $B_i=T(A_i)$ and $\bar B=\frac1N\sum_i B_i$. By Lemma~\ref{lem:mean_minimizer_app},
		$\frac1N\sum_i\normF{B_i-\bar B}^2\le \frac1N\sum_i\normF{B_i-T(\bar A)}^2$ with $\bar A=\frac1N\sum_i A_i$. Then $\normF{B_i-T(\bar A)}\le \normF{A_i-\bar A}$, and the claim follows.
	\end{proof}
	
	
	\begin{proof}[Proof of Lemma~\ref{lem:step-disp}]
		Fix $k$ and abbreviate $A_i:=A_{i,k}$, $A:=A_k^{\mathrm{ref}}$, $\tilde g_i:=\tilde g_{i,k}$, $\tilde{\bar g}:=\tilde{\bar g}_k$.
		Then $s_i=-A_i^{-1}\tilde g_i$ and $s^{\mathrm{ref}}=-A^{-1}\tilde{\bar g}$.
		Decompose:
		\[
		s_i-s^{\mathrm{ref}}
		=-A_i^{-1}(\tilde g_i-\tilde{\bar g}) + (A^{-1}-A_i^{-1})\tilde{\bar g}.
		\]
		Since $A_i=\tilde H_{i,k}+\tilde\lambda_{i,k}I\succeq \tilde\lambda_{i,k}I\succeq \underline\lambda_k I$,
		$\normtwo{A_i^{-1}}\le 1/\underline\lambda_k$.
		Thus
		\[
		\norm{s_i-s^{\mathrm{ref}}}
		\le \frac{1}{\underline\lambda_k}\norm{\tilde g_i-\tilde{\bar g}} + \normtwo{A^{-1}-A_i^{-1}}\,\norm{\tilde{\bar g}}.
		\]
		Use the resolvent identity:
		\[
		A^{-1}-A_i^{-1}=A_i^{-1}(A_i-A)A^{-1}.
		\]
		Therefore
		\[
		\normtwo{A^{-1}-A_i^{-1}}
		\le \normtwo{A_i^{-1}}\normtwo{A_i-A}\normtwo{A^{-1}}
		\le \frac{1}{\underline\lambda_k}\cdot \frac{1}{\tilde{\bar\lambda}_k}\,\normtwo{A_i-A},
		\]
		since $A\succeq \tilde{\bar\lambda}_k I$ implies $\normtwo{A^{-1}}\le 1/\tilde{\bar\lambda}_k$.
		Moreover
		\[
		A_i-A=(\tilde H_{i,k}-\tilde{\bar H}_k)+(\tilde\lambda_{i,k}-\tilde{\bar\lambda}_k)I,
		\]
		so $\normtwo{A_i-A}\le \normtwo{\tilde H_{i,k}-\tilde{\bar H}_k}+|\tilde\lambda_{i,k}-\tilde{\bar\lambda}_k|$.
		Combine and average over $i$, using
		$\norm{\bar s-s^{\mathrm{ref}}}\le \frac1N\sum_i \norm{s_i-s^{\mathrm{ref}}}$.
	\end{proof}
	
	\begin{proof}[Proof of Proposition~\ref{prop:inexact-rn}]
		The cubic regime $L_2\norm{\bar s_k}\le \tilde{\bar\lambda}_k$ is Lemma~\ref{lem:avg-regime}.
		For the residual, write
		\[
		r_k=(\nabla^2 f(\bar x_k)+\lambda_k I)\bar s_k+g_k,\qquad \lambda_k=\tilde{\bar\lambda}_k.
		\]
		Add and subtract $s_k^{\mathrm{ref}}$:
		\[
		r_k=(\nabla^2 f(\bar x_k)+\lambda_k I)(\bar s_k-s_k^{\mathrm{ref}})
		+(\nabla^2 f(\bar x_k)+\lambda_k I)s_k^{\mathrm{ref}}+g_k.
		\]
		Using $(\tilde{\bar H}_k+\lambda_k I)s_k^{\mathrm{ref}}=-\tilde{\bar g}_k$, we have
		\[
		(\nabla^2 f(\bar x_k)+\lambda_k I)s_k^{\mathrm{ref}}+g_k
		=(\nabla^2 f(\bar x_k)-\tilde{\bar H}_k)s_k^{\mathrm{ref}}+(g_k-\tilde{\bar g}_k).
		\]
		Thus
		\[
		r_k=(\nabla^2 f(\bar x_k)+\lambda_k I)(\bar s_k-s_k^{\mathrm{ref}})
		+(\nabla^2 f(\bar x_k)-\tilde{\bar H}_k)s_k^{\mathrm{ref}}+(g_k-\tilde{\bar g}_k).
		\]
		Take norms:
		\begin{align*}
			\norm{r_k}
			&\le \normtwo{\nabla^2 f(\bar x_k)+\lambda_k I}\,\norm{\bar s_k-s_k^{\mathrm{ref}}}
			+\normtwo{\nabla^2 f(\bar x_k)-\tilde{\bar H}_k}\,\norm{s_k^{\mathrm{ref}}}
			+\norm{g_k-\tilde{\bar g}_k}.
		\end{align*}
		First, $\normtwo{\nabla^2 f(\bar x_k)+\lambda_k I}\le M_H+\lambda_k$ by Assumption~\ref{ass:standing}(5).
		Using condition (1),
		\[
		\normtwo{\nabla^2 f(\bar x_k)+\lambda_k I}\,\norm{\bar s_k-s_k^{\mathrm{ref}}}
		\le (M_H+\lambda_k)\cdot \frac{\eta}{8}\cdot\frac{\lambda_k}{M_H+\lambda_k}\norm{\bar s_k}
		=\frac{\eta}{8}\lambda_k\norm{\bar s_k}.
		\]
		Second, by Lemma~\ref{lem:bridge-hess} and condition (3),
		$\normtwo{\nabla^2 f(\bar x_k)-\tilde{\bar H}_k}\le \eta\lambda_k/8$.
		Also
		\[
		\norm{s_k^{\mathrm{ref}}}\le \norm{\bar s_k}+\norm{\bar s_k-s_k^{\mathrm{ref}}}
		\le \norm{\bar s_k}+\frac{\eta}{8}\norm{\bar s_k}
		\le \frac98\norm{\bar s_k},
		\]
		since $\eta\le 1/12$ implies $1+\eta/8\le 9/8$.
		Hence
		\[
		\normtwo{\nabla^2 f(\bar x_k)-\tilde{\bar H}_k}\,\norm{s_k^{\mathrm{ref}}}
		\le \frac{\eta}{8}\lambda_k\cdot \frac98\norm{\bar s_k}=\frac{9\eta}{64}\lambda_k\norm{\bar s_k}.
		\]
		Third, by Lemma~\ref{lem:bridge-grad} and condition (2),
		$\norm{g_k-\tilde{\bar g}_k}\le \eta\lambda_k\norm{\bar s_k}/8$.
		Summing:
		\[
\norm{r_k}\le \left(\frac{\eta}{8}+\frac{9\eta}{64}+\frac{\eta}{8}\right)\lambda_k\norm{\bar s_k}
=\frac{25\eta}{64}\lambda_k\norm{\bar s_k}\le \eta\lambda_k\norm{\bar s_k}.
		\]
	\end{proof}
	
	\begin{proof}[Proof of Lemma~\ref{lem:lambda-comparable}]
		From \eqref{eq:rel-bridge}:
		\[
		\norm{g_k}\ge \norm{\tilde{\bar g}_k}-\norm{g_k-\tilde{\bar g}_k}\ge (1-\alpha_b)\norm{\tilde{\bar g}_k},
		\quad
		\norm{g_k}\le \norm{\tilde{\bar g}_k}+\norm{g_k-\tilde{\bar g}_k}\le (1+\alpha_b)\norm{\tilde{\bar g}_k}.
		\]
		Thus
		\[
		\frac{1}{1+\alpha_b}\norm{g_k}\le \norm{\tilde{\bar g}_k}\le \frac{1}{1-\alpha_b}\norm{g_k}.
		\]
		From \eqref{eq:rel-disp}, for all $i$,
		\[
		(1-\alpha_d)\norm{\tilde{\bar g}_k}\le \norm{\tilde g_{i,k}}\le (1+\alpha_d)\norm{\tilde{\bar g}_k}.
		\]
		Taking square roots and multiplying by $\sqrt{M}$:
		\[
		\sqrt{1-\alpha_d}\sqrt{M\norm{\tilde{\bar g}_k}}\le \lambda_{i,k}=\sqrt{M\norm{\tilde g_{i,k}}}
		\le \sqrt{1+\alpha_d}\sqrt{M\norm{\tilde{\bar g}_k}}.
		\]
		Averaging gives the same bounds for $\bar\lambda_k:=\frac1N\sum_i\lambda_{i,k}$.
		Since $\lambda_k=\bar\lambda_k+\bar e_k\ge \bar\lambda_k$,
		\[
		\lambda_k\ge \sqrt{1-\alpha_d}\sqrt{M\norm{\tilde{\bar g}_k}}
		\ge \frac{\sqrt{1-\alpha_d}}{\sqrt{1+\alpha_b}}\sqrt{M\norm{g_k}}.
		\]
		For the upper bound, \eqref{eq:e-absorb} implies
		$\lambda_k=\bar\lambda_k+\bar e_k\le \bar\lambda_k+\frac{\eta}{8}\lambda_k$,
		hence $(1-\eta/8)\lambda_k\le \bar\lambda_k$ and thus
		\[
		\lambda_k\le \frac{1}{1-\eta/8}\bar\lambda_k
		\le \frac{1}{1-\eta/8}\sqrt{1+\alpha_d}\sqrt{M\norm{\tilde{\bar g}_k}}
		\le \frac{1}{1-\eta/8}\cdot\frac{\sqrt{1+\alpha_d}}{\sqrt{1-\alpha_b}}\sqrt{M\norm{g_k}}.
		\]
	\end{proof}
	
	\begin{proof}[Proof of Lemma~\ref{lem:descent-inexact}]
		Apply Lemma~\ref{lem:taylor3} to $f$ at $(\bar x_k,\bar s_k)$:
		\[
		f(\bar x_k+\bar s_k)\le f(\bar x_k)+\ip{g_k}{\bar s_k}+\frac12\bar s_k^\top\nabla^2 f(\bar x_k)\bar s_k+\frac{L_2}{6}\norm{\bar s_k}^3.
		\]
		Let $d_k:=(\nabla^2 f(\bar x_k)+\lambda_k I)\bar s_k+g_k$.
		Then
		\[
		\ip{g_k}{\bar s_k}=-\bar s_k^\top\nabla^2 f(\bar x_k)\bar s_k-\lambda_k\norm{\bar s_k}^2+\ip{d_k}{\bar s_k}.
		\]
		Substitute:
		\[
		f(\bar x_{k+1})\le f(\bar x_k)-\frac12\bar s_k^\top\nabla^2 f(\bar x_k)\bar s_k-\lambda_k\norm{\bar s_k}^2+\ip{d_k}{\bar s_k}+\frac{L_2}{6}\norm{\bar s_k}^3.
		\]
		Since $\nabla^2 f(\bar x_k)\succeq 0$, drop the nonpositive term to get
		\[
		f(\bar x_{k+1})\le f(\bar x_k)-\lambda_k\norm{\bar s_k}^2+\ip{d_k}{\bar s_k}+\frac{L_2}{6}\norm{\bar s_k}^3.
		\]
		Bound $\ip{d_k}{\bar s_k}\le \norm{d_k}\norm{\bar s_k}\le \eta\lambda_k\norm{\bar s_k}^2$.
		Also $L_2\norm{\bar s_k}\le \lambda_k$ gives $\frac{L_2}{6}\norm{\bar s_k}^3\le \frac{1}{6}\lambda_k\norm{\bar s_k}^2$.
		Combine.
	\end{proof}
	
	\begin{proof}[Proof of Lemma~\ref{lem:grad-inexact}]
		By the integral form of Taylor's theorem and Lipschitz Hessian,
		\[
		\nabla f(\bar x_k+\bar s_k)=g_k+\nabla^2 f(\bar x_k)\bar s_k+r_k,\qquad \norm{r_k}\le \frac{L_2}{2}\norm{\bar s_k}^2.
		\]
		Using $d_k=(\nabla^2 f(\bar x_k)+\lambda_k I)\bar s_k+g_k$ we have
		$g_k+\nabla^2 f(\bar x_k)\bar s_k=-\lambda_k\bar s_k+d_k$.
		Thus
		\[
		\norm{g_{k+1}}\le \lambda_k\norm{\bar s_k}+\norm{d_k}+\norm{r_k}
		\le \lambda_k\norm{\bar s_k}+\eta\lambda_k\norm{\bar s_k}+\frac{L_2}{2}\norm{\bar s_k}^2.
		\]
		Finally $L_2\norm{\bar s_k}\le \lambda_k$ implies $\frac{L_2}{2}\norm{\bar s_k}^2\le \frac12\lambda_k\norm{\bar s_k}$.
	\end{proof}
	
	\begin{proof}[Proof of Lemma~\ref{lem:steady-rec}]
		Assume $\eta\le 1/12$. Then in Lemma~\ref{lem:descent-inexact},
		$c_d:=1-\eta-\frac16\ge \frac34$.
		Let $C_g:=1+\eta+\frac12\le 2$ (since $\eta\le 1/12$).
		For $k\in\mathcal I$, Lemma~\ref{lem:grad-inexact} yields
		\[
		\norm{g_{k+1}}\le C_g\lambda_k\norm{\bar s_k}.
		\]
		Because $k\in\mathcal I$, $\norm{g_{k+1}}\ge \frac14\norm{g_k}$, hence
		\[
		\norm{\bar s_k}\ge \frac{\norm{g_k}}{4C_g\lambda_k}.
		\]
		By Lemma~\ref{lem:descent-inexact},
		\[
		\Phi_k-\Phi_{k+1}\ge c_d\,\lambda_k\norm{\bar s_k}^2
		\ge c_d\,\lambda_k\left(\frac{\norm{g_k}}{4C_g\lambda_k}\right)^2
		=\frac{c_d}{16C_g^2}\frac{\norm{g_k}^2}{\lambda_k}.
		\]
		By Lemma~\ref{lem:lambda-comparable}, $\lambda_k\le C_\lambda\sqrt{M\norm{g_k}}$,
		thus
		\[
		\frac{\norm{g_k}^2}{\lambda_k}\ge \frac{1}{C_\lambda\sqrt{M}}\norm{g_k}^{3/2}.
		\]
		By convexity of $f$ and Assumption~\ref{ass:standing}(4),
		\[
		\Phi_k=f(\bar x_k)-f_\star \le \ip{g_k}{\bar x_k-x_\star}\le D\norm{g_k}
		\quad\Rightarrow\quad
		\norm{g_k}^{3/2}\ge D^{-3/2}\Phi_k^{3/2}.
		\]
		Combine to obtain $\Phi_k-\Phi_{k+1}\ge \tau \Phi_k^{3/2}$ with
		\[
		\tau:=\frac{c_d}{16C_g^2 C_\lambda D^{3/2}\sqrt{M}}.
		\]
	\end{proof}
	
	\begin{proof}[Proof of Lemma~\ref{lem:sharp-count}]
		Each $k\in\mathcal S$ satisfies $\norm{g_{k+1}}<\frac14\norm{g_k}$.
		After $m$ sharp steps, the gradient norm decreases by at least $4^{-m}$:
		$\norm{g_{k}}\le 4^{-m}\norm{g_0}$.
		To keep $\norm{g_k}>\varepsilon$ requires $4^{-m}\norm{g_0}>\varepsilon$, i.e., $m<\log_4(\norm{g_0}/\varepsilon)$.
	\end{proof}
	
	\begin{proof}[Proof of Lemma~\ref{lem:telescoping}]
		If there exists an index $k$ such that $u_k=\tau\Phi_k^{1/2}\ge 1$, then
		$\Phi_{k+1}\le \Phi_k-\tau\Phi_k^{3/2}=\Phi_k(1-u_k)\le 0$, hence $\Phi_{k+1}=0$
		and the claim holds trivially for all subsequent indices. Therefore, we may
		assume $u_k\in[0,1)$ for all $k$.
		Assume $\Phi_k>0$ and set $\psi_k:=\Phi_k^{-1/2}$.
		From $\Phi_{k+1}\le \Phi_k-\tau\Phi_k^{3/2}=\Phi_k(1-\tau\Phi_k^{1/2})$, define $u_k:=\tau\Phi_k^{1/2}\in[0,1)$ and get
		\[
		\psi_{k+1}=\Phi_{k+1}^{-1/2}\ge \Phi_k^{-1/2}(1-u_k)^{-1/2}=\psi_k(1-u_k)^{-1/2}.
		\]
		For $u\in[0,1)$, convexity of $(1-u)^{-1/2}$ implies $(1-u)^{-1/2}\ge 1+\frac{u}{2}$.
		Thus $\psi_{k+1}\ge \psi_k(1+u_k/2)=\psi_k+\tau/2$.
		Iterating: $\psi_k\ge \psi_0+\frac{\tau}{2}k\ge \frac{\tau}{2}(k+2)$, hence
		$\Phi_k=\psi_k^{-2}\le \frac{4}{\tau^2(k+2)^2}$.
	\end{proof}
	
	\begin{proof}[Proof of Theorem~\ref{thm:main}]
		We first show $\Phi_k=\cO(1/k^2)$.
		Let $m_k:=|\mathcal S\cap\{0,\dots,k-1\}|$.
		If $m_k\ge k/2$, then $\norm{g_k}\le 4^{-m_k}\norm{g_0}\le 4^{-k/2}\norm{g_0}$.
		By convexity and Assumption~\ref{ass:standing}(4), $\Phi_k\le D\norm{g_k}\le D\norm{g_0}4^{-k/2}$, which is $\cO(1/k^2)$.
		Otherwise $m_k<k/2$, so there are at least $k/2$ steady indices up to $k$.
Let $k_1<k_2<\cdots<k_m$ be the indices in $\mathcal I\cap\{0,1,\dots,k-1\}$,
so $m\ge k/2$. Define $\Psi_j:=\Phi_{k_j}$.
By Lemma~\ref{lem:descent-inexact}, $\{\Phi_\ell\}$ is nonincreasing, hence
$\Phi_{k_{j+1}}\le \Phi_{k_j+1}$ for all $j$.
Since $k_j\in\mathcal I$, Lemma~\ref{lem:steady-rec} gives
$\Phi_{k_j+1}\le \Phi_{k_j}-\tau\Phi_{k_j}^{3/2}=\Psi_j-\tau\Psi_j^{3/2}$.
Therefore, $\Psi_{j+1}\le \Psi_j-\tau\Psi_j^{3/2}$ for all $j$.
Applying Lemma~\ref{lem:telescoping} to $\{\Psi_j\}$ yields
$\Psi_m\le \frac{4}{\tau^2(m+2)^2}\le \frac{16}{\tau^2(k+4)^2}$,
and since $\Phi_k\le \Psi_m$, we obtain $\Phi_k=\cO(1/k^2)$.
		Applying Lemma~\ref{lem:telescoping} to that recursion yields $\Phi_k\le C/(k+2)^2$ for a constant $C$.
		Thus $\Phi_k=\cO(1/k^2)$ in all cases.
		
		Next, by Lemma~\ref{lem:gap-grad} with $h=f$ and $L=L_1$,
		$\norm{g_k}^2\le 2L_1\Phi_k$, hence $\norm{g_k}=\cO(1/k)$.
		Finally, $\norm{g_k}\le C/k\le \varepsilon$ holds when $k\ge C\varepsilon^{-1}$.
	\end{proof}
	
	\begin{proof}[Proof of Theorem~\ref{thm:comm}]
		Under \eqref{eq:log_schedule}, $\tau_k=t_k=\cO(\log(k+2))$. By Theorem~\ref{thm:end2end}, $K(\varepsilon)=\cO(\varepsilon^{-1})$. Hence
		\[
		\sum_{k=0}^{K(\varepsilon)-1}(\tau_k+t_k)\le C \sum_{k=0}^{K(\varepsilon)-1}\log(k+2)
		\le C K(\varepsilon)\log(K(\varepsilon)+2)=\cO(\varepsilon^{-1}\log(\varepsilon^{-1})).
		\]
	\end{proof}
	
\end{document}